\documentclass[11pt,a4paper]{article}

% =====================================================================
% [†ONT-TIME] CANONICAL TIME AND THE QUANTUM LOCK — this paper is the HOME (ONTOLOGY_FOUNDATION_INDEX.md §1n).
%   THE CORPUS'S ONE QUANTUM RESULT lives here (§lock), and p0's six-ways unification depends on it.
%   THE DIAGNOSIS: the frozen constraint is NOT a defect — "it is the FAITHFUL CANONICAL IMAGE OF THE BLOCK."
%   If the block is what exists, the formalism correctly reports there is nothing for a Hamiltonian to do.
%   THE NECESSITY, two grounds and ONE correction: occurrence is not existence (granting existence to the
%   manifold of occurrences "is a MISUSE OF THE WORD 'EXIST'"), and the CMB measures the cosmic rest frame.
%   "Declining the block and admitting the cosmic rest frame are THE SAME MOVE."
%   *** THE SELECTION IS THE WHOLE OF THE MOVE. The bare ADM machinery "has no preferred time of its own, and
%   in this exact sense the bare formalism contains NO RESOLUTION of the problem of time AND NO OBSTRUCTION to
%   one. Everything turns on what is fed to it." — "THE MOVE IS THE SELECTION. EVERYTHING ELSE IS READING THE
%   TEXTBOOK ON IT." We claim NO NOVELTY IN THE MECHANISM; what CR adds is a WARRANT. The MULTIPLE-CHOICE
%   PROBLEM has force "precisely when the choice of clock is a matter of formal convenience internal to the
%   theory"; it loses that force when the clock is fixed FROM OUTSIDE by ontology and observation — dissolved,
%   not solved within the space of formal clocks.
%   GUARD: the cosmic foliation is NOT one of the description groupoid's gauge orbits to be quotiented away —
%   "it is FIXED BY THE ONTOLOGY THE GAUGE FREEDOM IS DEFINED OVER." ***
%   Deparametrization: p_tau + H_phys = 0, solved for the generator => i d_tau Psi = H_phys Psi, UNITARY in
%   cosmic time. "Deparametrization is not a trick that restores a time that was really absent; it is the
%   formal face of a time THAT WAS NEVER LOST." The minisuperspace run is explicitly A TOY; and — a FEATURE,
%   not a gap — THE CR COSMOLOGY IS NON-SINGULAR, so "the usual motivation for quantizing minisuperspace, TO
%   RESOLVE AN INITIAL SINGULARITY, DOES NOT ARISE."
%   THE CLOSED-S^3 LIFT (the substantive canonical content): the TT graviton modes form a DISCRETE tower
%   (closed topology => discreteness), deparametrized to unitary evolution in cosmic time. TWO READINGS OF ONE
%   TOWER: on the ontological background a(T)=alpha.cosh(T/alpha); in a fundamental observer's frame
%   a(tau)=sinh^{2/3}, the discrete index passing to a continuous wavenumber. WHAT DISTINGUISHES THEM IS THE
%   BEND: the observer's Friedmann readout carries 2M/a^3, absent from the background's pure-Lambda expansion.
%   "The matter the observer attributes to the graviton's background is THE BEND OF THE SLICING — a feature of
%   the projection, NOT OF THE LAYER."
%   GUARD: "It is ONE CONGRUENCE, carrying the flat-LambdaCDM and the closed-S^3 slicings as TWO
%   SYNCHRONIZATIONS OF ITSELF, NOT TWO FRAMES."
%   LAMBDA>0 FORCES THE EXTERNAL CLOCK: an internal (area) clock is available at Lambda=0 and would serve,
%   "but at Lambda>0 it no longer separates the dynamics, and the absolute cosmic time is FORCED."
%   *** THE QUANTUM LOCK (§lock), AND IT CLOSES WITHOUT A FREE PARAMETER:
%   (i) ORDERING CANNOT DECIDE IT. In the geodesic coordinate x ~ a^{3/2} the kinetic operator carries an
%   inverse-square term gamma/x^2 whose coefficient, ACROSS THE NATURAL ORDERING FAMILY, attains a maximum
%   gamma = 1/4 — STRICTLY BELOW the essential-self-adjointness threshold 3/4 — so the origin is limit-circle
%   for EVERY ordering, while -(Lambda/8pi)a^3 maps to the marginal -x^2 and is limit-point at infinity.
%   DEFICIENCY INDICES (1,1), INDEPENDENTLY OF ORDERING.
%   (ii) THE BOUNDARY IS NOT A GENERIC ENDPOINT. a=0 IS THE DE SITTER COSMOLOGICAL HORIZON — the smooth,
%   finite-curvature metric singularity at which the comoving congruence's past null ruling closes. A Killing
%   horizon of surface gravity kappa carries a HARTLE-HAWKING thermal state (the unique Euclidean-regular
%   state, free of a conical defect); for the dS horizon kappa = 1/alpha, beta = 2.pi.alpha,
%   T = hbar/(2.pi.alpha.k_B) (Gibbons-Hawking).
%   (iii) EUCLIDEAN REGULARITY *IS* THE FRIEDRICHS EXTENSION: it retains the regular branch x^{1/2+nu} and
%   excludes the irregular admixture, "whose presence is a CONICAL DEFECT — an off-temperature state." So the
%   deficiency-(1,1) freedom is NOT MERELY PARAMETRIZED BUT CLOSED, AND CLOSED WITHOUT A FREE PARAMETER, since
%   the de Sitter horizon fixes ONE temperature. HBAR ENTERS THE GRAVITATIONAL SECTOR HERE, AT THE SEAM,
%   SCALED BY LAMBDA ALONE — the quantum instance of the substrate's constant ledger (p0, §1j).
%   (iv) THE COUPLED SECTOR, AND WHAT IS ACTUALLY OPEN. Deparametrizability is NOT AT RISK: the Brown-Kuchar
%   dust momentum enters the scalar constraint LINEARLY WHATEVER THE GRAVITON CONTENT, so the constraint solves
%   for a true Hamiltonian AT EVERY ORDER; the coupling enters H_phys as interaction terms, "not as an
%   obstruction." The boundary coefficient is promoted from the c-number gamma <= 1/4 to an OPERATOR Gamma-hat
%   straddling the 3/4 threshold; it decomposes as a DIRECT INTEGRAL — essentially self-adjoint where
%   Gamma-hat >= 3/4 (the graviton momentum makes the origin limit-point and REMOVES the boundary freedom),
%   limit-circle below. THE SAME THERMAL REGULARITY SUPPLIES THE CONDITION ON THE SUB-THRESHOLD FIBRES at the
%   one horizon period (kappa belongs to the BACKGROUND horizon, not the graviton content, so it is common to
%   every fibre). THEREFORE WHAT REMAINS OPEN IS NOT THE BOUNDARY CONDITION BUT THE *DEFINITION OF THE
%   INTERACTING TOWER* — the spectrum of Gamma-hat and the UV definition of the tower sums: "the standard
%   problem of the interacting theory rather than a residual freedom in the quantization at the boundary."
%   DO NOT REPORT THE COUPLED SECTOR AS A RESIDUAL QUANTIZATION FREEDOM. ***
%   DISSOLUTION, NOT SOLUTION: "a solution would add structure to make a defective formalism work; A
%   DISSOLUTION REMOVES A MISTAKEN PREMISE AND FINDS THE FORMALISM WAS WORKING ALL ALONG." ONE ONTOLOGICAL
%   CORRECTION, TWO CANONICAL PAYOFFS — the same correction that turns the frozen constraint into a true
%   Hamiltonian ALSO dissolves the inference that gravitational collapse completes in finite cosmic time.
%   P7 dissolves the problem of time at the AXIOMS ([†ONT-CR] §1e); this paper carries out the CONSTRUCTION.
%   P12 identifies it as the base-dependence of the substrate's coset metric ([†ONT-ALG] §1h): "the structural
%   and the canonical sides of ONE dissolution."
%
% =====================================================================
% CANON NOTE — LIVE OPERATING PHILOSOPHY
%
% CR is ONTOLOGICAL, not merely representational. The universe is a real
% 3-sphere evolving in a real de Sitter substrate; the layer S_t is what
% EXISTS and advances; the 4-manifold M is the representational record of
% what HAPPENS as it exists. "Representation/projection/shadow" mean
% built-by-construction-and-REAL, never unreal. The block universe is the
% error of granting the manifold of occurrences the existence that belongs
% only to the layer. Keep the ontological reading live throughout; do not
% let any phrasing collapse the layer's existence into "mere" representation.
%
% REGISTER: recognition, not proposal. The equations are flawless; only the
% reading shifts. The thesis is NOT "CR solves the problem of time" but
% "the problem of time was the block-universe category error in canonical
% clothing; here is the formalism read without it."
% =====================================================================

\usepackage[margin=1in]{geometry}
\usepackage{amsmath, amssymb, amsthm}
\usepackage{mathtools}
\usepackage{mathrsfs}
\usepackage{graphicx}
\usepackage[dvipsnames]{xcolor}
\usepackage{hyperref}
\usepackage{receipts}  % in-place receipt citations -> Appendix R
\usepackage{caption}

\theoremstyle{plain}
\newtheorem{theorem}{Theorem}
\newtheorem{corollary}{Corollary}
\newtheorem{lemma}{Lemma}
\newtheorem{proposition}{Proposition}
\theoremstyle{definition}
\newtheorem{axiom}{Axiom}
\newtheorem{definition}{Definition}
\theoremstyle{remark}
\newtheorem{remark}{Remark}

\newcommand{\St}{\mathcal{S}_t}
\newcommand{\Hphys}{H_{\mathrm{phys}}}
\newcommand{\dd}{\mathrm{d}}   % r2376+c54.9: used but undefined -- the paper did not compile

\title{The canonical problem of time as a category error:\\
\large an empirically forced cosmic time, deparametrization, the dissolution of frozen dynamics,\\
and a graviton quantization the de~Sitter horizon closes without a free parameter---the seam where $\hbar$ enters gravity---in Cosmological Relativity}
\author{Daryl Janzen\\
\small Department of Physics and Engineering Physics\\
\small University of Saskatchewan}
\date{}

\begin{document}
\maketitle

\begin{abstract}
In the canonical formulation of general relativity the Hamiltonian is a sum of constraints that vanish on the physical phase space: no preferred time survives, and the generator of evolution annihilates physical states rather than evolving them. This is the canonical ``problem of time.'' We argue that it is not a technical defect awaiting a technical repair, but the canonical symptom of a category error---treating the four-dimensional manifold of \emph{occurrences} as the thing that \emph{exists}, the block universe. The companion framework paper~\cite{JanzenCRframework} introduced, as an ontological commitment of Cosmological Relativity (CR), a global cosmic time that general relativity's formalism admits but does not supply. Here we show what that commitment does to the canonical formalism. The decisive premise is not a new calculation but a selection: the bare Arnowitt--Deser--Misner (ADM) machinery propagates data along an \emph{arbitrary} hypersurface and settles nothing, whereas CR selects a preferred cosmic foliation---forced empirically by the isotropy of the cosmic microwave background (CMB) and structurally by the horizon geometry of collapse~\cite{JanzenBHcausality}. On that foliation the lapse recovers its meaning as the metrical rate of the existing layer's advance, $d\tau = N\,dt$; the constraint, deparametrized along cosmic time, is solved for a \emph{true} Hamiltonian that generates the advance; and quantum evolution in cosmic time is unitary. The frozen constraint and the true Hamiltonian are the same canonical content under two readings, differing only in whether the manifold is granted existence. The deparametrization is, as a technique, entirely standard; what CR adds is not a mechanism but a \emph{warrant}. The multiple-choice problem---that many candidate internal times exist, that inequivalent quantizations follow from different choices, and that the formalism prefers none---has force exactly when the clock is chosen for formal convenience internal to the theory, and loses it when the clock is fixed from outside the formalism by ontology and observation. It is dissolved, not solved within the space of formal clocks. We exhibit the structure first in flat Friedmann--Lema\^itre minisuperspace, where the deparametrized Hamiltonian reproduces the flat-$\Lambda$CDM expansion---the $\sinh^{2/3}$ law established in~\cite{JanzenCRframework}---and then lift it to the closed-$S^3$ layer's propagating sector: the transverse-traceless graviton modes form a discrete tower that deparametrizes to a unitary evolution in cosmic time, and project into a fundamental observer's frame as the flat-$\Lambda$CDM graviton, with matter density appearing only as the bend of the slicing rather than as ontological content. The external clock is no convenience: an internal area clock separates the dynamics at $\Lambda=0$ and would serve, but at $\Lambda>0$ it no longer does, and the absolute cosmic time is forced. In that sector the quantization moreover closes \emph{without a free parameter}. The scale-factor Hamiltonian carries deficiency indices $(1,1)$ independently of operator ordering---the inverse-square coefficient at the origin attaining $\gamma=\tfrac14$ across the natural ordering family, strictly below the essential-self-adjointness threshold $\tfrac34$---so a single boundary condition at $a=0$ remains; and that point is not a generic half-line endpoint but the de~Sitter cosmological horizon, whose Hartle--Hawking thermal state, at the surface gravity $\kappa=1/\alpha$ the substrate itself fixes, selects the Friedrichs extension. The deficiency freedom is thereby not merely parametrized but closed, and closed without a free parameter: $\hbar$ enters the gravitational sector at the seam, scaled by $\Lambda$ alone. With the tower coupled, the boundary coefficient is promoted to an operator straddling the threshold and the same thermal regularity supplies the condition fibre by fibre, so that what remains open is not the boundary condition but the \emph{definition of the interacting tower}---the standard problem of the interacting theory rather than a residual freedom in the quantization. The problem of time is thereby \emph{dissolved} as a category error, not solved by an addition to the formalism, in exactly the way the same ontological correction dissolves the inference that gravitational collapse completes in finite cosmic time~\cite{JanzenBHcausality}.
\end{abstract}

\section{Introduction: the problem of time, as standardly posed}\label{sec:intro}

The canonical quantization of general relativity confronts a difficulty that has resisted resolution for half a century. In the ADM decomposition~\cite{ADM1962}, spacetime is foliated by spatial hypersurfaces and the Einstein--Hilbert action is recast in Hamiltonian form. The result is a fully constrained system in the sense of Dirac~\cite{Dirac1964}: the total Hamiltonian is a sum of the scalar (Hamiltonian) constraint $\mathcal{H}$ and the momentum constraint $\mathcal{H}_i$, each weighted by the lapse $N$ and shift $N^i$, and each vanishing on the physical phase space. The lapse and shift are not dynamical; they are Lagrange multipliers enforcing the constraints. Promoting the scalar constraint to an operator yields the Wheeler--DeWitt equation $\hat{\mathcal{H}}\Psi = 0$~\cite{DeWitt1967}: the physical state is annihilated by the generator that should evolve it. The wavefunction of the universe does not appear to depend on any time parameter at all.

This is the problem of time~\cite{Kuchar1992, Isham1993}. Its sharpest statement is that the formalism contains no preferred time variable, that every spatial foliation of a given spacetime is related to every other by a gauge transformation generated by the constraints, and that ``evolution'' in the coordinate label $t$ is therefore pure gauge---a re-slicing of one fixed four-dimensional object, carrying no physical content. The many proposed responses---identifying an internal time among the canonical variables before quantization, coupling a material clock, treating observables as relational ``evolving constants,'' or abandoning a fundamental time altogether---each carry well-known costs~\cite{Kuchar1992, Isham1993, Rovelli2004, BarbourEnd}, to which we return in Section~\ref{sec:relation}.

Our claim is that the difficulty is conceptual before it is technical, and that the conceptual diagnosis is precise rather than vague. The frozen formalism is the canonical expression of a single metaphysical commitment: that the four-dimensional manifold is the object that exists. We will argue, following the ontological foundation of CR~\cite{Janzen2025, JanzenCRframework} and the necessity assembled in Section~\ref{sec:necessity}, that this commitment is a category error---the confusion of what \emph{happens} with what \emph{exists}---and that it is empirically false, the universe having furnished, in the CMB, direct evidence of the very cosmic time the formalism declines to recognize. When that error is declined, the same canonical equations describe a real cosmic time all along.

We stress at the outset the register of this paper. We do not propose a modification of general relativity, a new term in the action, or an additional structure bolted onto the canonical formalism. The equations are those of Einstein and of ADM, unchanged. What changes is the reading. The frozen constraint and a true Hamiltonian generating cosmic-time evolution are, as we will show, the same content under two readings, differing only in whether the manifold of occurrences is granted the kind of existence that belongs to the evolving three-dimensional world. The problem of time is not solved here. It is dissolved.

The paper proceeds as follows. Section~\ref{sec:necessity} earns the premise the argument needs---that a preferred cosmic time is real---on two independent grounds, one conceptual and one empirical. Section~\ref{sec:two-readings} reads the unchanged ADM formalism under the two ontologies, and Section~\ref{sec:selection} isolates why the move is a \emph{selection} rather than a construction: the bare machinery settles nothing, and everything turns on whether a foliation is distinguished to be selected. Section~\ref{sec:deparam} carries out the deparametrization, structurally and then in flat minisuperspace, marking the boundary between the two. Section~\ref{sec:lock} lifts it to the closed-$S^3$ layer's propagating sector---the transverse-traceless graviton tower, advanced unitarily in cosmic time---and there the quantization closes: the scale-factor Hamiltonian's deficiency-$(1,1)$ freedom, ordering-independent, is fixed by the de~Sitter horizon's own thermal state, without a free parameter, leaving open not the boundary condition but the definition of the interacting tower. Section~\ref{sec:dissolution} states the verdict negatively, and Section~\ref{sec:relation} places the argument against the internal-time, relational, and timeless programmes, where the multiple-choice objection is met not with a mechanism but with a warrant.

\section{The necessity: a cosmic time that is real, not posited}\label{sec:necessity}

The canonical move of this paper rests on one premise: that there is a preferred cosmic time, and that it is a real feature of the world rather than a convention chosen for convenience. This premise is not assumed. It is forced, on two independent grounds---one conceptual, one empirical---which we summarize here; the full argument is developed in~\cite{Janzen2025, JanzenWildGoose, JanzenFortress} and stated as the necessary-and-sufficient augmentation of general relativity in~\cite{JanzenModernParallax}.

\paragraph{The conceptual ground: occurrence is not existence.}
To exist is to endure---to persist through an interval. The three-dimensional world exists in this sense, and so do the things in it. An \emph{event}, by contrast, is something that \emph{happens}: it occurs at a place and an instant and is gone. A worldline is the continuous succession of events associated with an enduring thing, and four-dimensional spacetime is the manifold of all such events---the totality of what happens. The persistent confusion about time, traceable from Heraclitus and Parmenides through to modern eternalism, is the treatment of occurrences as though they were existents: the reification of the event-manifold into a thing that exists~\cite{Janzen2025}. Once spacetime is imagined to exist, every moment along a worldline is imagined to exist on equal footing, the distinction between past, present, and future is declared an illusion, and time is said not to pass. But that conclusion was smuggled into the premise. Granting existence to the manifold of occurrences is not a discovery about time; it is a misuse of the word ``exist.'' The coherent picture is the one in which an evolving three-dimensional world exists and advances, and events happen in the course of that advance, mapping to a four-dimensional spacetime that represents them without being a further existent of the same kind.

In CR this distinction is structural. What exists is the layer $\St$---a real three-sphere evolving in a real de Sitter substrate~\cite{JanzenCRframework, JanzenSlicing}. The four-manifold $M$ is the representational record of the events that happen as the layer advances. The layer's existence and advance are not optional gloss on a fundamentally four-dimensional object; they are the ontology, and $M$ is its shadow.

\paragraph{The empirical ground: the CMB measures the cosmic rest frame.}
That a preferred cosmic time might be real but undetectable is the situation relativity leaves open: local experiments cannot distinguish a true state of rest, because the relativity of synchrony---a general kinematic consequence of one's chosen rest convention, independent of the light postulate---and the round-trip symmetry of local signalling together hide one's motion. Einstein's identification of simultaneity with frame-relative synchrony was, on this much, a philosophical choice among two empirically equivalent options, not a deduction; treating the local inaccessibility of a cosmic rest frame as evidence of its nonexistence is a modal fallacy~\cite{JanzenWildGoose}.

The universe, however, did furnish the missing evidence. The cosmological redshift of the CMB is isotropic to high precision, and the primordial temperature anisotropies at the level of one part in $10^5$ have been preserved across the entire history of expansion~\cite{Aghanim2020}. Had the expansion rate varied locally with the inhomogeneous matter distribution, the accumulated redshift along different lines of sight would differ by far more than is observed---by some three orders of magnitude, as we establish quantitatively in the companion paper~\cite{JanzenModernParallax}---and the isotropy therefore certifies that cosmological space has expanded uniformly throughout cosmic history. But to speak of a three-dimensional space expanding uniformly \emph{through time} already presupposes a cosmic time, a cosmic rest frame, and an associated simultaneity. The same data fix our own velocity relative to that frame through the CMB dipole. The cosmic rest frame is thus not a large-scale average or a coordinate convenience but one of the better-measured structures in cosmology. Synchrony is relative; simultaneity is absolute~\cite{JanzenWildGoose, JanzenFortress}.

\paragraph{The two grounds are one correction.}
These are not two arguments but two faces of one. The relativity-of-simultaneity premise is precisely what licenses the reification of the block: if there is no objective ``now,'' the four-dimensional manifold is all there coherently is, and the layer's privileged advance has no referent. Conversely, an empirically real cosmic ``now'' restores the layer as the existent and demotes the manifold to its record. Declining the block (the conceptual correction) and admitting the cosmic rest frame (the empirical correction) are the same move. The equations of relativity are untouched by it; only the interpretation shifts---from ``if two events appear simultaneous they must be simultaneous'' to ``apparent simultaneity is a local matter with no bearing on the global cosmic simultaneity that cosmology establishes empirically''~\cite{JanzenFortress}.

This is the premise the canonical argument needs, and it has been earned rather than assumed: there is a preferred cosmic foliation, and it is real.

\section{The canonical formalism, read two ways}\label{sec:two-readings}

We now read the standard canonical formalism under the two ontologies. Nothing in this section alters the equations.

In the ADM form, a spacetime $(M,g)$ is foliated by spacelike hypersurfaces $\Sigma_t$ with induced metric $h_{ij}$ and conjugate momentum $\pi^{ij}$, and the line element is
\begin{equation}
ds^2 = -N^2\,dt^2 + h_{ij}\,(dx^i + N^i\,dt)(dx^j + N^j\,dt),
\end{equation}
with lapse $N$ and shift $N^i$. The action becomes $S = \int dt\,d^3x\,\big(\pi^{ij}\dot h_{ij} - N\mathcal{H} - N^i\mathcal{H}_i\big)$, and variation of $N$ and $N^i$ enforces the constraints $\mathcal{H}=0$, $\mathcal{H}_i=0$. The total Hamiltonian $H = \int d^3x\,(N\mathcal{H} + N^i\mathcal{H}_i)$ is a sum of constraints; on the constraint surface it vanishes.

\paragraph{The frozen reading.}
Suppose the existing object is the four-dimensional manifold itself. Then a choice of foliation is a choice of how to slice one already-existing thing into a stack of three-dimensional pictures, and a different foliation is a different slicing of the same thing. No slicing is distinguished; the relation between any two is a gauge transformation generated by the constraints; and since the generator of $t$-evolution is a constraint, it moves a physical state nowhere. Quantum-mechanically, $\hat{\mathcal{H}}\Psi=0$: the state is timeless. This reading is internally consistent. Its content is exactly the content of the premise: \emph{if} the block is what exists, time does not pass, and the formalism correctly reports that there is nothing for a Hamiltonian to do. The frozen formalism is not a mistake of calculation. It is the faithful canonical image of the block.

\paragraph{The reading CR adopts.}
Suppose instead, with CR, that what exists is the evolving layer $\St$, and that $M$ is the record of the events that happen as it advances (Section~\ref{sec:necessity}). Then a foliation is not a slicing of a pre-existing block; it is a choice of how to coordinate the advance of the existing world, and one such foliation---the cosmic one---tracks that advance rather than cutting across it. On this reading the lapse is not a bookkeeping multiplier but the metrical amount by which the existing layer advances per unit of the time label,
\begin{equation}
d\tau = N\,dt,
\end{equation}
with the shift $N^i$ carrying the foliation's non-orthogonality to the layer. The constraint $\mathcal{H}=0$ is still an equation; but now it is read as the relation that, once solved for the momentum conjugate to the cosmic time, yields the generator of the layer's advance---a true Hamiltonian. The same symbols, the same constraint surface, the same dynamics; a different answer to the single question of what exists. This lapse and shift are the very levels the cosmological development computes on: the lapse is the foliation stacking rate the observable expansion rides---radiation-free, the density the content that set rate carries, read off the clock rather than sourcing it---and the shift is the synchronization through which that expansion is \emph{observed} as distance and redshift, with the leaf's bend the matter it carries~\cite{JanzenOperator, JanzenCosmogenesis}. That a plasma process rides the one rate or the other is then fixed by which side of the reassignment seam it falls on, not chosen: the same selection of the existent's foliation that turns the frozen constraint into a true Hamiltonian here fixes, one level down, which rate a recombination-era or a nucleosynthesis-era computation must read---and reifying the shift's synchronous slicing as the existent is the block error, met at the level of a rate.

\section{The selection is the whole of the move}\label{sec:selection}

It is worth isolating why the foregoing is a recognition rather than a new construction, because this is the crux of the paper.

The ADM machinery propagates canonical data along whatever foliation one supplies. It is indifferent to which foliation that is; it has no preferred time of its own, and in this exact sense the bare formalism contains no resolution of the problem of time and no obstruction to one. Everything turns on what is fed to it. If nothing distinguishes a foliation, the constraint is all there is and the state is frozen. If a foliation is distinguished, the constraint deparametrizes along it and a true Hamiltonian appears. The canonical formalism is the same in both cases; what differs is whether a preferred foliation exists to be selected.

CR supplies exactly that, and supplies it from outside the bare formalism, on the grounds of Section~\ref{sec:necessity}: the cosmic foliation is real, forced empirically by the CMB and conceptually by the occurrence/existence distinction. It is moreover the foliation that the rest of the programme independently singles out. The metric-singularity structure of the event horizon determines a unique limiting causal orientation in collapse, generically non-orthogonal to any spacelike slice, and this is the orientation CR reassigns as cosmic time~\cite{JanzenBHcausality, JanzenCRframework}; the representational freedom in choosing among Lorentzian metrics on the fixed manifold is a genuine gauge symmetry, organized by the description groupoid~\cite{JanzenGroupoid}, but the cosmic foliation is not one of its orbits to be quotiented away---it is fixed by the ontology the gauge freedom is defined over.

The consequence is methodological as well as physical. There is no separate ``canonical machinery'' result to be proved here; the deparametrization that follows is standard once a clock is in hand. The content of this paper is that CR's cosmic time is the physically correct clock, on grounds that are ontological and empirical rather than formal, and that the canonical face of the existing physics is therefore a true Hamiltonian rather than a frozen constraint. The move is the selection. Everything else is reading the textbook on it.

\emph{The same selection settles what a mass charge could be.} The ADM mass is defined relative to an asymptotic time translation, and exists only where the geometry supplies one---which asymptotically-de~Sitter spacetimes do not, there being no global timelike Killing vector, so that no conserved charge is well defined there at all~\cite{JanzenSlicing}. \emph{On the reading taken here that absence is not a deficiency of the geometry but a misplacement of the question}: the time this construction runs on is not recovered from an asymptotic symmetry but selected and \emph{measured}, so a quantity whose definition waits on an asymptotic Killing vector is waiting on the wrong thing. It is the same move as the paper's own---what the formalism cannot supply from inside, the selection supplies from outside, and the question that presupposed the missing structure lapses with it.

\section{Deparametrization and the true Hamiltonian}\label{sec:deparam}

We now carry out the reading. The structural statement is independent of any model; the explicit minisuperspace realization that follows is an illustration of it, and we are careful to mark the boundary between the two.

\paragraph{The structural statement.}
Let the cosmic time of Section~\ref{sec:necessity} be the distinguished foliation parameter, and let $p_\tau$ denote the momentum conjugate to it. The scalar constraint, linear in $p_\tau$ after the cosmic clock is isolated, takes the deparametrized form
\begin{equation}
p_\tau + \Hphys\big(q^A, p_A\big) = 0,
\label{eq:deparam}
\end{equation}
where $(q^A,p_A)$ are the remaining (true, gauge-invariant) degrees of freedom of the layer and $\Hphys$ depends on the cosmic time only through them. Equation~\eqref{eq:deparam} is the same constraint that, read with no preferred foliation, gives $\hat{\mathcal{H}}\Psi=0$. Read with the cosmic clock selected, it is solved for the generator of advance: classically $dq^A/d\tau = \partial \Hphys/\partial p_A$, and on quantization
\begin{equation}
i\,\partial_\tau \Psi(q^A,\tau) = \hat{\Hphys}\,\Psi(q^A,\tau).
\label{eq:schrodinger}
\end{equation}
This is a genuine time-dependent Schr\"odinger equation generating unitary evolution in cosmic time, not the frozen Wheeler--DeWitt constraint. The foliation is fixed---by the ontology, not by a gauge choice---so there is a true time and a true Hamiltonian. The frozen constraint and the true Hamiltonian are \eqref{eq:deparam} read two ways; they are the same content, differing only in whether the manifold is granted existence. Deparametrization here is not a trick that restores a time that was really absent; it is the formal face of a time that was never lost, because the cosmic foliation was real all along.

\paragraph{A zero-duration boundary, and what it does to the evolution.} The cosmic time of
Section~\ref{sec:necessity} is the real part of the companion framework's complex cosmic time, and that
framework's closed contour contains one segment---the \emph{lift}, joining the collapse's comoving turnaround
to the branch point at which a collapse becomes a cosmology---along which \emph{the real part does not
advance}~\cite{JanzenCRframework}. The consequence for \eqref{eq:schrodinger} is immediate and worth stating
plainly, because it is unusual: cosmogenesis occupies \emph{no cosmic time at all}, so the evolution operator
across it is $U(\Delta\tau=0)=\mathbb{1}$. The true Hamiltonian has no interval in which to act there. The
geometry nonetheless changes across the segment---the areal radius climbs from the turnaround to zero and the
expansion rate is carried from zero to divergent---so the change is not generated by $\Hphys$ but is the
segment's own analytic content, the contour read at another point of itself. \emph{The beginning is therefore
not a first instant of the evolution this section constructs; it is a boundary of zero duration within it}, and
that is why nothing can be lost across it: no evolution acts, so none can be non-unitary.

We mark the limit of the claim. What is established is the classical statement---the foliation parameter does
not advance on that segment, so no interval of $\tau$ separates the two legs---together with its immediate
consequence for the evolution operator. \emph{What is not established here is a quantization of the Euclidean
segment itself.} The state on either side is related by the identity in cosmic time; whether the imaginary-time
stretch admits a quantum treatment in its own right, and what that treatment would say about the
$\{q^A,p_A\}$ across it, is not addressed by this deparametrization and is not claimed by it. \emph{The gap is named rather than managed}: it is carried as an open problem of the companion framework~\cite{JanzenCRframework}, and its shape is definite. \emph{The silence is structural, not a matter of refinement}---an evolution parametrized by cosmic time cannot describe a segment of zero cosmic duration, so no sharpening of $\hat{\Hphys}$ reaches it. What is owed is an account of an imaginary-time stretch across which the classical configuration changes while the foliation parameter does not, and it is the natural point of contact between this construction and Euclidean approaches to the quantum beginning---contact the programme does not presently make. \emph{The classical account of the beginning is complete; the quantum account of it is absent}, and the second half of that sentence travels with the first. One step of it can nonetheless be taken here, and taking it sharpens what remains.

\paragraph{The propagator across the lift is Euclidean.} The evolution \eqref{eq:schrodinger} is generated in
the framework's complex cosmic time, and on the lift that time runs along the imaginary direction:
$\tilde\tau=\mathrm{const}+is$, so $\partial_s=i\,\partial_{\tilde\tau}$ and \eqref{eq:schrodinger} becomes
$\partial_s\Psi=-\hat{\Hphys}\Psi$. The kernel carrying the state across is therefore not the unitary
$U=e^{-i\hat{\Hphys}\Delta\tau}$ but
\begin{equation}
K=e^{-\hat{\Hphys}\,\lvert\Delta\eta\rvert},
\label{eq:euclidean-kernel}
\end{equation}
a Euclidean kernel---real, contracting, and not unitary---with $\lvert\Delta\eta\rvert$ the lift's imaginary
conformal-time interval. This is consistent with $U(\Delta\tau=0)=\mathbb{1}$ rather than in tension with it:
\emph{no cosmic time elapses, and the kernel is not an evolution in cosmic time.}

\emph{Two things follow, and the first is a check the construction did not have to pass.} A mode of frequency
$\omega$ is damped by $e^{-\omega\lvert\Delta\eta\rvert}$ under \eqref{eq:euclidean-kernel}; the framework's
independent classical reading of the same segment damps a mode carrying $e^{ikc_s\eta}$ by
$e^{-kc_s\lvert\Delta\eta\rvert}$~\cite{JanzenCRframework}. With $\omega=kc_s$ these are the same expression,
term for term. \emph{The classical selection rule---frozen content crosses, oscillating content does
not---is the Euclidean kernel, read classically.} Second, and this is the content the classical reading could
not supply: a Euclidean kernel run for a long imaginary interval \emph{projects onto the ground state},
suppressing each excitation by $e^{-(E_n-E_0)\lvert\Delta\eta\rvert}$. \emph{The scope of that statement should be fixed carefully, because the two sectors it separates are not symmetric.} The kernel acts on \emph{oscillatory} content; a frozen, zero-frequency mode is a fixed point of it and passes unchanged. And at the crossing no mode is oscillating: on the contracting leg $aH$ grows without bound as $r\to0$, so the comoving horizon shrinks to zero and every mode exits it and freezes before the branch point~\cite{JanzenCRcosmology}. \emph{So the projection does not select the frozen content that crosses---it has nothing to act on there.} What it acts on is the oscillatory content of modes while they were sub-horizon, which it annihilates; the ground-state projection is the statement that this content, and not the frozen amplitude, is what the segment removes.

\emph{The projection is not exact, and where it fails is a scale rather than a caveat.} At
$\lvert\Delta\eta\rvert\simeq1.7\times10^{4}\,\mathrm{Mpc}$ the damping reaches $e^{-2}$ only at
$k\sim2\times10^{-4}\,\mathrm{Mpc}^{-1}$, so modes below that wavenumber are not strongly suppressed and cross
with residual excitation. That wavenumber subtends $\ell\approx3$ on the last-scattering sphere, which is
where the framework's independent transmission calculation places the boundary of the filter's
reach~\cite{JanzenCRframework}. The two routes agree on the scale as well as on the factor.

\emph{The obvious objection is that a Euclidean kernel needs a Hamiltonian bounded below, and that Euclidean
gravity notoriously lacks one---the conformal factor entering the Einstein--Hilbert action with the opposite
sign, so that the Euclidean action can be driven arbitrarily negative. \emph{That objection does not reach this
construction, and the reason is structural rather than fortunate.}} The conformal-factor problem arises when the
path integral ranges over the conformal factor of the metric. Here it does not. The substrate's scale is
$\alpha$ and is fixed---not chosen, but required, as the unique maximally symmetric structure carrying no
unforced modulus~\cite{JanzenShadowExistence, JanzenSlicing}---so there is no conformal mode to integrate over;
and the propagating gravitational degree of freedom of the layer is the \emph{transverse-traceless} shear of
Section~\ref{sec:lock}, which is precisely the sector the conformal mode is absent from. Mode by mode that
sector is \eqref{eq:tt-action}, a harmonic oscillator, whose Hamiltonian is bounded below. The areal radius
along the lift is likewise confined, running between the comoving turnaround and the branch point on an
interval fixed by the progenitor mass, with the substrate curvature finite throughout~\cite{JanzenCRframework,
JanzenCircle}: there is no runaway direction for the kernel to diverge along.

\emph{And the state the kernel projects onto is not a new posit but one this paper has already selected on
independent grounds.} Section~\ref{sec:lock} finds the scale-factor Hamiltonian to have deficiency indices
$(1,1)$---a one-parameter family of candidate quantizations---and closes it, without stipulation, by observing
that the boundary at $a=0$ is the de~Sitter horizon, that its surface gravity is $\kappa=1/\alpha$, and that the
regular Euclidean state at that $\kappa$ leaves nothing to choose. \emph{That is the same condition the kernel
\eqref{eq:euclidean-kernel} enforces}: regularity in the Euclidean continuation. \emph{And that regularity is available only because of what the causality companion establishes about the boundary.} A regular Euclidean state exists at a horizon only if the Euclidean continuation is smooth there---free of a conical defect, curvature finite. The horizon at $a=0$ is a \emph{metric} singularity, at which the spatial measure collapses while the curvature stays finite, and not a curvature singularity~\cite{JanzenBHcausality}; had it been the latter there would be no smooth Euclidean section, no regular state to select, and the extension would remain unfixed. \emph{The distinction that dissolves the information paradox is the same distinction that makes this quantization unique.} The extension-fixing condition
and the lift's kernel are one requirement met twice, once at the horizon and once across the beginning.

\emph{What remains open should be stated at its true size, which is smaller than it was.} The kernel's frequency
is not constant along the segment---the layer's scale varies with the path parameter---so the projection is
adiabatic rather than exact, and the residual at the largest scales quantified above is one face of that. And
the relation between the configuration variables $\{q^A,p_A\}$ on the two sides is described here only through
the kernel's action on the modes, not derived from a variational principle for the segment itself. \emph{The gravitational sector is not the obstruction, and the segment's variational treatment can be given
here.}

\paragraph{The lift is an instanton, and its action is finite.} The contour's own law is a first integral: for
the marginal congruence $(\dd r/\dd\tilde\tau)^{2}+f=1$, which is $\tfrac12\dot r^{2}+V=\tfrac12$ with
$V(r)=\tfrac12 f(r)$. Continuing $\tilde\tau=is$ sends $\dot r\mapsto-i\,r'$ and the Lorentzian action
$S=\int\dd\tilde\tau\,[\tfrac12\dot r^{2}-V]$ to $S=-iS_{E}$ with
\begin{equation}
S_{E}=\int\dd s\;\Big[\tfrac12 r'^{2}+V(r)\Big],
\label{eq:euclidean-action}
\end{equation}
whose stationary points obey $r''=+V'(r)$---motion in the \emph{inverted} potential, the standard signature of
a Euclidean solution---with first integral $\tfrac12 r'^{2}-V=-\tfrac12$. The lift's closed form
$r(s)=-(2M\alpha^{2})^{1/3}\lvert\sin(3s/2\alpha)\rvert^{2/3}$ satisfies both: the first integral is
$-\tfrac12$ identically along the segment, and $r''$ agrees with $V'(r)$ to the precision of the
check\rcpt{LIFT_instanton_action}. \emph{So the imaginary-time segment is not a formal device but a solution of
a variational principle}---an instanton connecting the Lorentzian turning point to the branch point.

Its action converges. Near the branch point $V\simeq-M/r$ while $r\propto s^{2/3}$, so the integrand of
\eqref{eq:euclidean-action} grows as $s^{-2/3}$ and is integrable; numerically
$S_{E}\to1.489$ in units $\alpha=1$ on the forced member. \emph{The segment therefore carries a finite weight
rather than a divergent one}, which is what a saddle-point treatment of the beginning would require of it.
\emph{Two normalisations are in play, and the relation between them can be settled rather than flagged.} The
functional \eqref{eq:euclidean-action} carries a unit kinetic term; the gravitational action carries the metric
factor and the prefactor implied by \eqref{eq:Hphys}. They share a solution---the contour's law
$(\dd r/\dd\tilde\tau)^{2}=2M/r+r^{2}/\alpha^{2}$ is the Friedmann equation with $a=r$, the identification
$2M\leftrightarrow(8\pi/3)\rho_{d}$ being just $M=(4\pi/3)\rho_{d}$---so the instanton is one trajectory read in
two measures, and only its \emph{action value} differs.

That value follows without further input. On shell the reduced action is $S=\int p_{a}\,\dd a$, and inverting
\eqref{eq:Hphys} gives $p_{a}=(3/4\pi)a\dot a$; continuing $\tilde\tau=is$ yields $S=-iS_{E}$ with
\begin{equation}
S_{E}^{\mathrm{grav}}=\frac{3}{4\pi}\int\dd s\;r\,r'^{2}
=\frac{3}{4\pi}\int\dd s\;r\,\bigl[f(r)-1\bigr],
\label{eq:grav-action}
\end{equation}
using $r'^{2}=f-1$ on the segment. The integral converges rapidly and gives
$S_{E}^{\mathrm{grav}}=-0.0481\,\alpha^{2}/G$ on the forced member\rcpt{LIFT_gravitational_action}.

\emph{Its sign is the informative part, and it is not a pathology.} A negative Euclidean gravitational action is
the Hartle--Hawking sign: the no-boundary de~Sitter action is likewise negative, $-\alpha^{2}/8G$ in this
convention, yielding an enhanced rather than suppressed weight~\cite{HartleHawking1976}. The present value is
the same sign and the same order, smaller by a factor of order two, which is what a segment of the contour
rather than a full hemisphere should give. The sign traces to the segment lying on the $r<0$ branch, where
$p_{a}<0$ while $r$ increases. \emph{So the beginning's weight is of the no-boundary type, obtained here from a
contour the construction already possessed rather than from a boundary condition imposed on the path integral.} \emph{One register point must be made here, because the vocabulary invites a reading the construction
refuses.} The segment's parametrisation runs imaginary, and the kernel and action computed above are the
kernel and action of that parametrisation; \emph{they are not the Euclidean objects of a Wick-rotated
spacetime}. The geometric-core paper states the guard generally: the construction reaches a real manifold
through imaginary instruments, and the continuations it uses---the embedding coordinate, the equatorial seam,
the cosmogenesis reassignment---are \emph{real analytic continuations} on a spacetime that is Lorentzian
throughout, not Wick rotations~\cite{JanzenGeometricCore}. The de~Sitter horizon's Gibbons--Hawking state, by
contrast, \emph{is} a Euclidean continuation, and that paper marks it as distinct in kind from these. \emph{So
the two must not be conflated}: the thermal register in which $\hbar$ is fixed by the period
$\beta=2\pi\alpha$ is not the register the segment's action belongs to, and dividing that action by that
$\hbar$ would be joining two continuations the construction keeps apart. What the action is, is a real integral
along a real curve of a real Lorentzian geometry, whose parameter happens to run imaginary over one stretch;
what it is not is an exponent awaiting a quantum of action. \emph{Its comparison with the no-boundary sign
above is a comparison of signs and magnitudes, not an identification of frameworks.}
\paragraph{The adiabatic correction, and where it fails.} The kernel projects exactly only if
$\hat{\Hphys}$ is constant along the segment, and it is not: the tensor tower of Section~\ref{sec:lock} carries
frequencies $\omega_{n}=\mu_{n}/a$, and $a$ varies along the lift, diverging in $\omega$ as the branch point is
approached. The exact suppression exponent is therefore $\int\omega_{n}\,\dd s$ rather than
$\omega_{n}\lvert\Delta\eta\rvert$, and the first question is whether it converges at all. It does: near the
branch point $\lvert r\rvert\propto s^{2/3}$, so the integrand grows only as $s^{-2/3}$ and
\begin{equation}
\int_{0}^{\pi\alpha/3}\frac{\dd s}{\lvert r(s)\rvert}=3.3387\,\alpha^{-1}
\label{eq:adiabatic-exponent}
\end{equation}
on the forced member\rcpt{LIFT_adiabatic_correction}. \emph{The correction is not small and it runs the
conservative way}: the constant-frequency estimate using the turnaround value gives $1.4396\,\alpha^{-1}$, so
the exact exponent is larger by a factor $2.32$ and the suppression is \emph{stronger} than the naive reading
of \eqref{eq:euclidean-kernel} suggests.

\emph{Whether the projection is adiabatic is a separate question with a definite answer.} The adiabaticity
parameter $\lvert\dd\omega/\dd s\rvert/\omega^{2}$ is $C/\mu_{n}$ with $C$ running from $1.72$ near the branch
point to $0.16$ near the turnaround, so the approximation is controlled by the harmonic index alone. With
$\mu_{n}^{2}=n(n+2)-2$, $n\ge2$, this gives $0.70$ at $n=2$, $0.48$ at $n=3$, and $0.16$ by $n=10$.
\emph{The projection is therefore adiabatic for all but the lowest few harmonics, and degrades to order unity
only at $n=2$ and $n=3$.} That the approximation should fail exactly where the tower is coarsest is expected
rather than surprising---there are no modes below $n=2$ on $S^{3}$---and it locates the correction rather than
leaving it as a caveat. \emph{We record, without claiming it, that the harmonic indices at which the treatment
loses control are the lowest ones}; the relation between $S^{3}$ tensor harmonics and observed multipoles is
not developed here, and no inference to the microwave sky is drawn from this section.

\paragraph{An explicit illustration, and its limits.}
To show the structure is not empty we exhibit it in the simplest cosmological model, flat Friedmann--Lema\^itre minisuperspace. With scale factor $a$, conjugate momentum $p_a$, cosmological constant $\Lambda$, and a pressureless (dust) reference fluid serving as the cosmic clock in the manner of Brown and Kucha\v{r}~\cite{BrownKuchar1995}, the deparametrized gravitational Hamiltonian reduced from the Einstein--Hilbert action (in units $G=1$, fiducial comoving volume unity) is
\begin{equation}
\Hphys(a,p_a) = \frac{2\pi}{3}\,\frac{p_a^2}{a} - \frac{\Lambda}{8\pi}\,a^3,
\label{eq:Hphys}   % r2376+c54.9: a second \label{eq:minisuper} stood here; amsmath keeps only one
%   per equation and silently dropped it, leaving two \eqref{eq:Hphys} dangling.  Both names pointed at
%   this same equation, so the references were redirected to eq:Hphys rather than the label duplicated.
\end{equation}
equal on shell to the conserved comoving dust energy\rcpt{P10_minisuperspace_friedmann}. Hamilton's equation $\dot a = \partial \Hphys/\partial p_a = 4\pi p_a/3a$, together with $\Hphys = \rho_d$ on shell, gives
\begin{equation}
\left(\frac{\dot a}{a}\right)^2 = \frac{8\pi}{3}\,\frac{\rho_d}{a^3} + \frac{\Lambda}{3},
\label{eq:friedmann}
\end{equation}
the flat-$\Lambda$CDM Friedmann equation. Its solution is the $\sinh^{2/3}$ law
\begin{equation}
a(\tau) = A\,\sinh^{2/3}\!\Big(\tfrac{3\tau}{2\alpha}\Big),\qquad \alpha = \sqrt{3/\Lambda},
\end{equation}
with the single scale $\alpha$ of the de Sitter substrate; one verifies directly that this $a(\tau)$ satisfies~\eqref{eq:friedmann} identically. This is precisely the expansion history that the companion framework paper derives, by a wholly independent route, as the SdS expansion law of the layered construction~\cite{JanzenCRframework}: the deparametrized generator reproduces, in cosmic time, the same flat-$\Lambda$CDM evolution that the null-boundary correspondence delivers geometrically. Promoting~\eqref{eq:Hphys} to an operator on $L^2$ of the half-line and applying~\eqref{eq:schrodinger} gives unitary evolution of the cosmological wavefunction in cosmic time.

Two honest qualifications attach to this illustration. First, it is a minisuperspace toy: a homogeneous truncation with a dust clock, not the full Schwarzschild--de Sitter object of CR, and the operator ordering of $p_a^2/a$ and the self-adjoint extension on the half-line are genuine technical questions not settled by the reduction. The illustration establishes that the deparametrized structure is consistent and reproduces the right cosmology; it does not by itself constitute the canonical quantization of the full theory. That next object---the lift to the closed-$S^3$ layer and its propagating degrees of freedom---is carried out in Section~\ref{sec:lock}. Second---and this is a feature rather than a gap---the CR cosmology is non-singular: the $\sinh^{2/3}$ law has no finite-cosmic-time curvature singularity, and indeed no finite layer reaches one~\cite{JanzenBHcausality, JanzenCRframework}. The usual motivation for quantizing minisuperspace, to resolve an initial singularity, therefore does not arise. The point of~\eqref{eq:schrodinger} is not singularity avoidance but the demonstration that, on the selected foliation, the canonical generator evolves rather than annihilates.

\section{The closed-$S^3$ lift: the graviton sector on the layer}\label{sec:lock}

The minisuperspace illustration of Section~\ref{sec:deparam} truncates the layer to its scale factor alone; it carries no propagating degree of freedom and so cannot exhibit what the true Hamiltonian generates beyond the background expansion. The full layer does carry such a degree of freedom---the transverse-traceless shear of its spatial geometry, the graviton---and the lift of the deparametrization to it is the substantive canonical content. We carry out that lift here. Throughout, the foliation parameter is the same absolute cosmic time of Section~\ref{sec:necessity}: the comoving congruence of the layer, which the companion framework paper identifies with the cosmic rest frame the CMB measures and with the $S^3$ family of the null-boundary correspondence~\cite{JanzenCRframework}. It is one congruence, carrying the flat-$\Lambda$CDM and the closed-$S^3$ slicings as two synchronizations of itself, not two frames.

The layer is best represented by the background geometry on which the cosmology is built---the reassigned de~Sitter geometry, distinct (as a Lorentzian metric) from the Schwarzschild--de~Sitter cosmology it underlies but sharing its manifold and foliation~\cite{JanzenCRframework}. Its closed synchronous slicing is the evolving round three-sphere of radius $a(T)=\alpha\cosh(T/\alpha)$ in cosmic time $T$. A transverse-traceless perturbation of the spatial metric decomposes into the tensor harmonics of $S^3$, $h_{ij}(T,x)=\sum_n \phi_n(T)\,Y^{(n)}_{ij}(x)$, with the $Y^{(n)}_{ij}$ the standard transverse-traceless rank-two harmonics of the unit three-sphere and Laplace eigenvalues $\mu_n^2=n(n+2)-2$, $n\ge 2$. The second-order Einstein--Hilbert action in the transverse-traceless sector reduces, mode by mode, to that of a harmonic oscillator with time-dependent mass $a^3$ and frequency $\mu_n/a$,
\begin{equation}
S_n=\tfrac12\!\int\! dT\,a^3\Big[\dot\phi_n^2-\frac{\mu_n^2}{a^2}\,\phi_n^2\Big],
\qquad \pi_n=a^3\dot\phi_n .
\label{eq:tt-action}
\end{equation}
Read on the absolute foliation exactly as in Section~\ref{sec:deparam}, the constraint deparametrizes and \eqref{eq:schrodinger} becomes
\begin{equation}
i\,\partial_T\Psi=\hat\Hphys\,\Psi,\qquad
\hat\Hphys=\sum_n\Big[\frac{\hat\pi_n^2}{2\,a(T)^3}+\tfrac12\, a(T)\,\mu_n^2\,\hat\phi_n^2\Big],
\label{eq:lock-schrodinger}
\end{equation}
a unitary evolution in cosmic time of a tower of time-dependent oscillators\rcpt{P10_graviton_lift}, one per tensor harmonic. This is the lift the toy lacked: the abstract gauge-invariant degrees of freedom $(q^A,p_A)$ of Section~\ref{sec:deparam} are here concrete---the graviton modes $\phi_n$---and the deparametrized generator advances them rather than annihilating them. Each mode is a Schr\"odinger oscillator on $L^2(\mathbb{R})$ and is manifestly self-adjoint; the closed topology of the layer enters as the \emph{discreteness} of the tower, in contrast to the continuous spectra of the flat minisuperspace and planar reductions.

The same tower has a second reading, in the frame of a fundamental observer. Under the null-boundary reassignment the observer's worldline is one of the reassigned null rulings, and the layer's expansion read along it is the areal radius $a(\tau)=(2M\alpha^2)^{1/3}\sinh^{2/3}(3\tau/2\alpha)$ in the observer's proper time $\tau$---the flat-$\Lambda$CDM scale factor of the framework paper, the expansion leg of the \emph{cosmogenetic bead} the framework proves one closed curve~\cite{JanzenCRframework}. The cosmic time deparametrized here is the very parameter along that bead: real on the expansion leg the observer inhabits, it is what continues off the real axis, through a bounded contour, along the bead's lap into the antecedent collapse. The tower projected into this frame is \eqref{eq:lock-schrodinger} with $a(T)$ replaced by $a(\tau)$, the discrete index passing to the continuous wavenumber of the observer's locally flat slices. What distinguishes the two readings is the bend: the observer's Friedmann readout is
\begin{equation}
\Big(\frac{\dot a}{a}\Big)^2=\frac{2M}{a^3}+\frac{1}{\alpha^2},
\label{eq:obs-friedmann}
\end{equation}
carrying a pressureless matter term $2M/a^3$ that is absent from the background geometry's pure-$\Lambda$ expansion, $(\dot a/a)^2\!\to\!\alpha^{-2}$. The matter the observer attributes to the graviton's background is thus the bend of the slicing---a feature of the projection, not of the layer---exactly the reclassification the framework makes for matter content generally~\cite{JanzenCRframework,JanzenOperator}. The graviton sector therefore exhibits, in one object, both faces of the layered reading: the discrete unitary tower on the ontological background, and the continuous flat-$\Lambda$CDM spectrum a real observer measures, with the matter density appearing only in the latter.

This realizes concretely the loop drawn in Section~\ref{sec:dissolution} with the dynamical content. The propagating degree of freedom is the transverse-traceless shear of the evolving layer, advanced by a true Hamiltonian on the absolute foliation---the same structure the companion dynamics paper exhibits for the confined polarized wave~\cite{JanzenDynamics}, here for the cosmological sector and the full tensor tower. The necessity of the external clock is sharpest here: an internal (area) clock is available at $\Lambda=0$ and would serve, but at $\Lambda>0$ it no longer separates the dynamics, and the absolute cosmic time is forced. Of the two background-sector questions, the first is settled in structure. The operator ordering and the half-line self-adjoint extension of the scale factor (Section~\ref{sec:deparam}) are one coupled question: in the geodesic coordinate $x\propto a^{3/2}$ the kinetic operator carries an inverse-square term $\gamma/x^2$ at the origin whose coefficient, across the natural ordering family, attains a maximum $\gamma=\tfrac14$---strictly below the essential-self-adjointness threshold $\gamma=\tfrac34$---so the operator is limit-circle at $a=0$ for every ordering, while the $-\tfrac{\Lambda}{8\pi}a^3$ term maps to the marginal $-x^2$ and is limit-point at $a\to\infty$, forcing no condition there. The scale-factor Hamiltonian therefore has deficiency indices $(1,1)$ independently of the ordering~\cite{Weyl1910, ReedSimon1975}: a one-parameter family of self-adjoint extensions fixed by a single boundary condition at $a=0$. The ordering ambiguity cannot render the quantization unique; the residual reduces to the physical choice of that boundary condition. That choice is fixed, for the free scale-factor sector, by the geometry of the boundary itself.
It is worth naming what kind of closure that is, because the corpus has a criterion for exactly this and this is its sharpest instance. A one-parameter family of self-adjoint extensions is a one-parameter family of \emph{candidate quantizations}---distinct physical predictions indexed by a $\lambda$ that the formalism does not pin. That is the structure the epistemic reading excludes: a structure carrying an unforced parameter is not a single world but a family, and answers ``what is the world?'' with a family rather than a world~\cite{JanzenShadowExistence}. What follows below is not a stipulation that removes the family but the substrate \emph{requiring} its own configuration: the boundary at $a=0$ is the de~Sitter horizon, its surface gravity is $\kappa=1/\alpha$, and the regular Euclidean state at that $\kappa$ leaves nothing to choose. \emph{The distinction that makes this the criterion's instance and not merely an analogy is worth keeping: $\lambda$ indexes distinct worlds, so it is an ontological family and inadmissible as it stands, whereas a parameter that merely records an unknown datum of one world---the progenitor's mass, say---is an epistemic gap and not a family at all.} The point $a=0$ is not a generic half-line endpoint but the de~Sitter cosmological horizon---the smooth, finite-curvature metric singularity at which the comoving congruence's past null ruling closes, the substrate curvature staying finite while the comoving ruler collapses~\cite{JanzenBHcausality, JanzenOperator, JanzenCRframework}. A Killing horizon of surface gravity $\kappa$ carries a thermal (Hartle--Hawking) state~\cite{HartleHawking1976}, the unique state regular in the Euclidean continuation and free of a conical defect at the horizon; for the de~Sitter cosmological horizon $f=1-r^2/\alpha^2$ gives $\kappa=\tfrac12|f'(\alpha)|=1/\alpha$, regularity period $\beta=2\pi/\kappa=2\pi\alpha$, and the Gibbons--Hawking temperature $T=\hbar\kappa/2\pi k_B=\hbar/2\pi\alpha k_B$~\cite{GibbonsHawking1977}. Imposed on the scale-factor wavefunction, Euclidean regularity is exactly the condition that retains the regular indicial branch $x^{1/2+\nu}$, $\nu=\sqrt{\gamma+\tfrac14}<1$, and excludes the irregular admixture $x^{1/2-\nu}$, whose presence is a conical defect---an off-temperature state. This is the Friedrichs extension, and it is a single boundary condition: the deficiency-$(1,1)$ freedom is thereby not merely parametrized but \emph{closed}, and closed without a free parameter, since the de~Sitter horizon fixes one temperature. The physical choice the ordering ambiguity left open is the horizon's own thermal state, and $\hbar$ enters the gravitational sector here, at the seam, scaled by $\Lambda$ alone---the quantum instance of the substrate's constant ledger, in which the fundamental constants enter as unit gauges over the single scale and the one place a free quantum parameter could sit is closed by the horizon's own thermal state~\cite{JanzenGeometricCore}. \emph{And that continuation is not a device brought in to close the freedom; it is the substrate's other real form.} The global Wick rotation $x_{0}\mapsto ix_{0}$ carries $\mathrm{dS}_{5}=SO(5,1)/SO(4,1)$ to the compact $S^{5}=SO(6)/SO(5)$, and the two are the real forms of the one complex group $SO(6,\mathbb{C})$~\cite{JanzenBoundary,JanzenGeometricCore}. \emph{So the register in which $\hbar$ enters is the same substrate read on its compact face}, and the horizon---where the period is $\beta=2\pi\alpha$---is where the two faces meet: the extension is fixed not by an extra assumption but by the geometry the Lorentzian reading already carries, seen from its other side. This determines the free scale-factor sector; the coupled sector below is where the open structure remains. The second background-sector question---the couplings of the tower to the scale-factor sector beyond quadratic order---sharpens this structure rather than threatening it. Deparametrizability is not at risk: the cosmic clock is a Brown--Kucha\v{r} dust fluid (Section~\ref{sec:deparam}), whose momentum enters the scalar constraint linearly whatever the graviton content, so the constraint solves for a true Hamiltonian at every order of the coupling---which enters $\Hphys$ as interaction terms between $a$ and the modes, not as an obstruction. The free tower above evolves on $a(T)$ as a fixed classical background, which is why each mode is manifestly self-adjoint on $L^2(\mathbb{R})$; the coupling question is what happens once the scale factor is itself quantized and back-reacts, and there it acts at the $a=0$ boundary of the scale-factor half-line. Each mode's kinetic term $\pi_n^2/2a^3$ is, in the geodesic coordinate, the inverse-square operator $\pi_n^2/2x^2$ at the origin, so the boundary coefficient is promoted from the c-number $\gamma\le\tfrac14$ of the free scale factor to an operator $\hat\Gamma$ on the tower whose spectrum straddles the $\tfrac34$ threshold. \emph{At leading order that spectrum is bounded below by $\gamma$ and unbounded above---the operator is a sum of squares, displayed below---but the decomposition that follows uses only that both sides of the threshold are occupied, and not that the spectrum has a floor; whether the complete $\hat\Gamma$ is bounded below is not assumed here~\rcpt{P10_the_straddle_does_not_need_a_floor}, though it does in fact follow: the full inverse-square coefficient is positive on non-degenerate metrics, so $\hat\Gamma=\gamma+c\sum_n\hat\pi_n^2$ retains its floor at $\gamma$ beyond leading order, and the question of a spectrum below $-\tfrac14$---where $\nu=\sqrt{\Gamma+\tfrac14}$ would leave the reals and no regular branch would exist---has no object. What remains open at this paragraph's end is accordingly not the floor but the straddle itself as a computed fact---whether the spectrum does occupy both sides of $\tfrac34$, which the decomposition assumes and which is a different threshold from $-\tfrac14$, deciding a different property: $\tfrac34$ separates limit-point from limit-circle and so decides whether a boundary condition must be chosen at all, while $-\tfrac14$ decides whether a regular branch exists to choose. How the straddle falls does not bear on the closure below, which is supplied fibre by fibre and so cannot be broken by the size of the sub-threshold set~\rcpt{D1_the_boundary_is_per_fibre_and_the_UV_is_over_fibres}.} As $\hat\Gamma$ commutes with the radial part, $-\partial_x^2+\hat\Gamma/x^2$ decomposes as a direct integral over its spectrum: essentially self-adjoint where $\hat\Gamma\ge\tfrac34$---the graviton momentum makes the origin limit-point and removes the boundary freedom---and limit-circle where $\hat\Gamma<\tfrac34$. The single boundary condition of the free scale factor is thereby replaced, in the coupled theory, by a self-adjoint condition supported on the sub-threshold graviton subspace, its deficiency that subspace rather than a line. At leading order $\hat\Gamma=\gamma+c\sum_n\hat\pi_n^2$; but the cubic and higher self-interactions enter at the same inverse-square order at the origin ($\pi_n^2\phi_m/a^3$ in kind), so the complete boundary coefficient is the singular part of the interacting graviton Hamiltonian, and its spectrum together with the ultraviolet definition of the tower sums is the open frontier---the interacting tower as a defined theory, met here at its boundary face. \emph{Two things can be said about that spectrum without settling the frontier, and both bear on the decomposition. First, it is computable branch by branch: on the natural ordering family the minimum of $\hat\Gamma$ is $\tfrac14$ under normal ordering and $\tfrac34$ under symmetric ordering with one mode occupied, and the difference between the two is exactly the zero-point quantum of that mode---which is why the threshold sits at $\tfrac34=\tfrac14+\tfrac12$, the free boundary coefficient plus one such quantum, rather than at a value the construction chose. Second, the two orderings answer the paper's own question oppositely: under normal ordering the origin remains limit-circle and the boundary freedom survives, while under symmetric ordering it becomes limit-point from the first occupied mode upward. The decomposition below is untouched by the choice, since it uses only that both sides of the threshold are occupied and both orderings occupy both sides; but the physical content is not, and the choice is not one this construction can make internally. Asking which ordering is asking whether the graviton tower's zero-point energy gravitates at the horizon---the cosmological-constant problem in local dress, reached from inside the boundary coefficient rather than imported.} The boundary condition on that sub-threshold subspace, however, is supplied by the same principle that fixed the free sector. Thermal (Hartle--Hawking) regularity is Euclidean smoothness at the horizon mode by mode: across the direct integral it imposes the regular branch $x^{1/2+\nu}$, $\nu=\sqrt{\hat\Gamma+\tfrac14}$, on each sub-threshold fibre at the one horizon period $\beta=2\pi\alpha$ (the surface gravity $\kappa=1/\alpha$ belongs to the background horizon, not to the graviton content, and so is common to every fibre), and asks nothing of the limit-point fibres $\hat\Gamma\ge\tfrac34$. The self-adjoint condition on the sub-threshold subspace is therefore the de~Sitter horizon's thermal state at every order of the coupling, exactly as in the free sector; what remains open is accordingly not the boundary condition but the \emph{definition of the interacting tower}---the spectrum of $\hat\Gamma$ and the ultraviolet definition of the tower sums---the first of which admits a verdict here, and it is affirmative\rcpt{P10_gamma_hat_is_bounded_below}. \emph{The DeWitt form carries exactly one negative direction and it is the conformal one}, which is why the Wheeler--DeWitt equation is hyperbolic; on transverse-traceless momenta that term drops and the form is $\mathrm{tr}((g\pi)^{2})$, equal to $\mathrm{tr}(S^{2})$ for the symmetric $S=L^{\mathsf T}\pi L$ with $L$ the Cholesky factor of $g$, hence positive definite for every non-degenerate spatial metric. \emph{So the cubic term's apparent unboundedness is an artefact of truncation}: $\pi^{2}(1-\lambda\phi+\cdots)$ is the expansion of $\pi^{2}/(1+\lambda\phi)$, whose full coefficient is positive wherever the metric is non-degenerate.\rcpt{D50_the_floor_does_survive_and_the_paper_said_so_232_revisions_earlier} \emph{That sentence is the one a reader arriving at the truncated form will not have reached yet, and the distance is short enough to be worth marking}: the expansion is handed over four sentences earlier, and a truncation of a positive function reads as an operator that has lost its floor. \emph{The cited receipt exists because that reading was taken}---the first-order term $-\lambda\phi\pi^{2}$ sampled alone does go arbitrarily negative, and the resummed coefficient on the same points never falls below $\gamma$. \emph{The deparametrization that yields a true Hamiltonian is the same move that yields a bounded-below one, both being the removal of the conformal direction}---though the positivity speaks only of the interior, the degenerate boundary being what the thermal condition above is for. What remains genuinely open is accordingly the ultraviolet definition of the tower sums---the standard problem of the interacting theory rather than a residual freedom in the quantization at the boundary. \emph{That is a characterisation, and it can be replaced by a
measurement and a structural remark, neither of which was in print.}\rcpt{S50_the_counterterm_basis_is_one_dimensional_because_the_background_family_is}
\emph{The measurement first.} The tower is a discrete sum rather than a momentum integral---the
harmonics run from $n\ge2$ with no zero mode and no soft region---so the infrared is regulated by the
compactness of the layer and needs nothing. The ultraviolet is untouched by that: with
$\mu_n\sim n$, $\langle\hat\pi_n^2\rangle\sim n$ in the instantaneous ground state, and a degeneracy
growing as $n^2$, the shell contribution goes as $n^3$ and the sum \emph{diverges quartically}---the
ordinary zero-point degree of a field in four dimensions. \emph{The third of those inputs was carried
as a scaling and is now a formula}: since $S^{3}=SU(2)$ is parallelizable, a frame-indexed field of
frame-spin $s$ is $L^{2}(SU(2))\otimes V_{s}$, Peter--Weyl returns level-$j$ totals of $1$, $3$ and $5$
times $(2j+1)^{2}$ for $s=0,1,2$---the component counts---and the transverse-traceless part is the two
extreme summands, exactly two fifths of the symmetric-tracefree total. Organised by eigenvalue,
\emph{the degeneracy is $2(n-1)(n+3)$}, ten at the floor $n=2$, so the shell contribution is $2n^{3}$
and the leading constant is settled rather than
assumed\rcpt{D1_the_degeneracy_carrying_the_quartic_was_never_derived_and_its_constant_is_the_component_count}.
\emph{And the constant is the propagating-component count rather than a universal}: the scalar tower's
counting function grows with leading coefficient $\tfrac13$ against the tensor tower's $\tfrac23$, so
the growth is $n^{2}$ per propagating component and not per tensor rank. \emph{So ``standard'' is the right word and
it is now a number.} \emph{And the leading term's counterterm is one this framework already carries.}
The quartic piece of the one-loop vacuum energy is independent of mass and of field---it survives at
$m=0$ and is the same for every species---so it is a constant vacuum energy, and \emph{the counterterm
a constant vacuum energy requires is a cosmological-constant term}. That is the framework's single
dimensionful constant, absorbed into the one observed curvature with no bare-$\Lambda$-versus-vacuum
split~\cite{JanzenGeometricCore}; and since $\ell_P$ is a gauge-combination rather than a second
physical length, the cutoff is not smuggling a scale in either.

\emph{The structural remark is about the terms after it, and it is sharper than the leading case
because their counterterms are not cosmological-constant terms.} Only the \emph{quadratic and
logarithmic successors} carry the mass, which is to say only they renormalise masses and couplings; in
the standard curved-space expansion the first goes with the Ricci scalar and the second with
curvature-squared invariants, and a curvature-squared coupling is not an entry in this framework's
ledger. \emph{The reason that is not fatal here is a property of the backgrounds this construction
admits rather than of the divergences.} On a maximally symmetric geometry of radius $\alpha$ every
curvature invariant is a pure power of $1/\alpha^{2}$~\cite{JanzenGeometricCore}, so
$\int\!\sqrt{g}\,R^{2}$, $\int\!\sqrt{g}\,R_{\mu\nu}R^{\mu\nu}$,
$\int\!\sqrt{g}\,R_{\mu\nu\rho\sigma}R^{\mu\nu\rho\sigma}$ and $\int\!\sqrt{g}$ are not four
independent functionals but four multiples of one; and the substrates the construction admits are a
\emph{one-parameter} family, the $\alpha$-foliation, every descent step above the last changing the
scale and nothing else~\cite{JanzenGeometricCore, JanzenAlgebroid}. \emph{So the counterterm basis is
one-dimensional because the admitted background family is}, and a divergence of any degree requires
exactly one counterterm where a generic theory requires three.

\emph{What that argument does not reach is stated here rather than left for a reader to find, because
it is the whole of what remains.} A degeneracy among counterterms is a property of the background
class: the functionals are distinguished by their response to \emph{varying} the background, and on a
one-parameter family they cannot be told apart. \emph{The first thing to check is therefore which
background the tower is defined on, and it is this section's own}: the closed synchronous slicing
$a(T)=\alpha\cosh(T/\alpha)$ of \S\ref{sec:deparam}, whose Ricci scalar is the constant $12/\alpha^{2}$.
\emph{That geometry is exactly de~Sitter, so the degeneracy holds on the very background the free tower
uses}, and the argument above is not an appeal to a nearby object.

\emph{Where it genuinely stops is elsewhere, and the section has already named the place.} The free
tower evolves on $a(T)$ as a \emph{fixed classical} background; the coupled sector is what happens once
the scale factor is itself quantized and back-reacts, and that is precisely the regime the boundary
operator $\hat\Gamma$ belongs to. \emph{The natural worry there is that a counterterm basis is a
statement about a class of fixed backgrounds, so that once the scale factor is quantized there is no
fixed background left to state it on. That worry does not survive being checked, and what replaces it
is sharper.} The collapse of the three quadratic invariants is not a consequence of maximal symmetry:
in four dimensions their deficit $\mathrm{Riem}^{2}-2\,\mathrm{Ric}^{2}+\tfrac13R^{2}$ is exactly
$C_{\mu\nu\rho\sigma}C^{\mu\nu\rho\sigma}$, every Friedmann geometry is conformally flat for
\emph{every} scale factor, and on such a geometry the Gauss--Bonnet combination is an exact total
derivative, $\sqrt{g}\,(R^{2}-4\,\mathrm{Ric}^{2}+\mathrm{Riem}^{2})=24\,\dd(\dot a^{3}/3+\dot
a)/\dd T$. So $\int\!\sqrt g\,\mathrm{Ric}^{2}$ and $\int\!\sqrt g\,\mathrm{Riem}^{2}$ are both
fixed by $\int\!\sqrt g\,R^{2}$ up to a boundary term whatever $a$ does: \emph{no scale factor breaks
it, because no scale factor can make an FRW geometry anything but conformally
flat}\rcpt{Q1_the_degeneracy_is_conformal_flatness_not_maximal_symmetry_so_no_scale_factor_can_break_it}.
\emph{And because that is an identity holding pointwise in $a(\cdot)$ rather than an evaluation on a
chosen class, it survives superposition and descends to the quantized sector as an operator relation}:
the coupled sector never required a class of fixed backgrounds, because the statement was never made by
evaluating on one. \emph{What maximal symmetry did buy is a weaker and different thing, and that half
does not survive}: on a constant-curvature background terms of \emph{different} dimension---$\int\!
\sqrt g$, $\int\!\sqrt g\,R$, $\int\!\sqrt g\,R^{2}$---are proportional as well, and they part
company as soon as $a$ is not the de~Sitter $\cosh$, as the radiation-free layer's running Ricci scalar
shows directly. \emph{So the place the degeneracy genuinely ends is not the scale factor but the shear.}
The deficit \emph{is} $C^{2}$; an anisotropic shear of amplitude $\sigma$ over an isotropic expansion
gives $C^{2}=4\sigma^{2}+O(\sigma^{4})$; and the propagating content of this tower is the
transverse-traceless shear. \emph{So what remains is the tower's own shear, which is a calculation and
not a question about meaning}---entering at second order in the mode amplitude, with the sub-leading
heat-kernel coefficients as its instrument. \emph{That calculation is run, and its answer is smaller
than the question sounds}: for a transverse-traceless perturbation the first variation of the Ricci
scalar vanishes identically and $\sqrt g$ is unchanged, so $R^{2}$ at second order in the amplitude is
$2\bar R$ times $R$ at that order---\emph{pointwise, with no integration by parts}---and
$\int\!\sqrt g\,R^{2}$ is accordingly a multiple of the Einstein--Hilbert term rather than a new
structure. With the Gauss--Bonnet combination contributing no field equation, the three quadratic
invariants leave $\int\!\sqrt g\,R^{2}$ and $\int\!\sqrt g\,C^{2}$, of which only the second is
new: \emph{the shear therefore costs exactly one new counterterm, and it is the Weyl-squared
one}\rcpt{S1_the_shear_needs_exactly_one_new_counterterm_and_my_own_count_was_one_too_many}. \emph{One
caveat belongs in the statement rather than after it}: the dimension-four scalars also include the
parity-odd Pontryagin density, which is not a curvature-squared invariant of the kind counted above
and is a total derivative, but which is non-zero at second order for a \emph{circularly} polarised
mode and reverses with the handedness---\emph{a linearly polarised mode returns zero for it, and the
corpus carries a chirality}~\cite{JanzenDynamics}, so the zero is a property of the mode chosen and
not of the geometry. The propagating sector itself is a clean, unitary, deparametrized evolution---a true Hamiltonian generating the graviton on the layer, in cosmic time.

\section{Dissolution, not solution}\label{sec:dissolution}

The result of Sections~\ref{sec:selection}--\ref{sec:deparam} is best stated negatively. We have added nothing to the formalism. We have not introduced a new time variable into general relativity, nor a matter field whose role is to be a clock, nor a modification of the constraint algebra. We have identified, on external grounds, which foliation is physically real, and read the existing constraint on it. The problem of time, on this account, was never a problem internal to the formalism; it was the formalism faithfully reporting the consequences of an ontological premise---the block---that is both a category error and empirically false.

This is why the appropriate verdict is dissolution rather than solution, and the distinction is not merely rhetorical. A solution would add structure to make a defective formalism work. A dissolution removes a mistaken premise and finds the formalism was working all along. The four epistemic rules that govern such corrections~\cite{JanzenFortress, JanzenShadowExistence}---resist the naïve identification of appearance with reality; prefer a reading on which the observed structure is a structural necessity rather than a tuned permission; consolidate phenomena under one principle; treat accumulating patches as a red flag---each favour the CR reading over the frozen one. The cosmic time is a structural necessity of the ontology and an empirical datum, not a tuned addition. The reading consolidates: the same correction that turns the frozen constraint into a true Hamiltonian also dissolves the inference that gravitational collapse completes in finite cosmic time, since both inferences---the timeless wavefunction and the completed horizon---rest on granting existence to a manifold of occurrences that no observer's present ever contains~\cite{JanzenBHcausality, JanzenFortress}. One ontological correction, two canonical payoffs.

The reading also closes a loop with the dynamical content of the programme. The true Hamiltonian generates the advance of the layer $\St$; the geometry of that advance is what the operator and range papers read directly, where the matter content of a layer is the bend of its spatial cut off the vacuum profile and gravitational radiation is the transverse-traceless shear of the evolving layer, up to the wall at which a symmetric layer can no longer carry the wave and ordinary inhomogeneous evolution begins~\cite{JanzenOperator, JanzenRange}. The canonical face (a Hamiltonian that advances the layer) and the geometric face (matter and radiation as features of the advancing layer's intrinsic geometry) are two descriptions of the one evolving three-dimensional world. The cosmogenesis re-expansion is one such advance: the companion cosmogenesis paper~\cite{JanzenCosmogenesis} runs the collapse-and-re-expansion thermal history---the layer's advance along cosmic time, generated by this true Hamiltonian rather than frozen---through a nuclear network to the primordial light-element abundances. The companion foundation roots these faces in a single augmentation of general relativity~\cite{JanzenModernParallax}; the present paper supplies the canonical one, and with it the recognition that CR's distinctive contribution to quantum gravity is not a new dynamics but an absolute cosmic time---real, forced, and already implicit in the formalism---that turns the constraint into a generator. This dissolution is the canonical face of the substrate's maximal symmetry gathered in the geometric core~\cite{JanzenGeometricCore}: the problem of time is identified in the constraint-algebroid construction as the base-dependence of the substrate's coset metric~\cite{JanzenAlgebroid}, and the selection of the forced foliation is what deparametrizes it away---the structural and the canonical sides of one dissolution.

\section{Relation to existing approaches}\label{sec:relation}

The deparametrization carried out in Section~\ref{sec:deparam} is, as a technique, entirely standard, and it is important to be clear about what is and is not new here. Kucha\v{r}'s survey~\cite{Kuchar1992} classifies responses to the problem of time into those that locate an internal time before quantization, those that quantize first and seek time afterwards, and those that abandon time outright; Isham's review~\cite{Isham1993} maps the same terrain. The construction of Section~\ref{sec:deparam} sits squarely in the first class, and its minisuperspace realization uses a dust clock of exactly the Brown--Kucha\v{r} type~\cite{BrownKuchar1995}. We claim no novelty in the mechanism.

What CR adds is not a mechanism but a warrant. The standard and well-known objection to internal-time approaches is the multiple-choice problem: many candidate internal times exist, inequivalent quantizations follow from different choices, and the formalism provides no principle to prefer one~\cite{Kuchar1992, Isham1993}. This objection has force precisely when the choice of clock is a matter of formal convenience internal to the theory. It loses that force when the clock is fixed from outside the formalism by ontology and observation. CR's cosmic time is not selected because it makes the equations tractable; it is selected because the occurrence/existence distinction identifies the evolving layer as what exists, and because the CMB measures the corresponding cosmic rest frame (Section~\ref{sec:necessity}). The multiple-choice problem is thus not solved within the space of formal clocks but dissolved by the recognition that one clock is physically real and the others are not. This is the same pattern as the rest of the paper: a difficulty that is genuine so long as a mistaken premise is held, and that disappears with the premise.

The timeless and relational programmes---Page and Wootters' conditional probabilities~\cite{PageWootters1983}, Rovelli's relational quantum mechanics and evolving constants of motion~\cite{Rovelli2004}, Barbour's denial of fundamental time~\cite{BarbourEnd}---take the opposite path with full consistency: they accept that no preferred foliation exists and reconstruct the appearance of time from correlations within a fundamentally timeless structure. We do not claim these programmes are internally flawed. We claim that the premise they share---that there is no objective cosmic now---is the very premise the CMB refutes~\cite{JanzenModernParallax}, and that once a real cosmic time is admitted, the elaborate reconstruction of temporal appearance from timeless correlations is solving a problem that no longer arises. The choice between the programmes is, in the end, the choice of Section~\ref{sec:necessity}: whether to grant the manifold of occurrences the existence that the evidence assigns to the evolving layer.

\section{Conclusion}\label{sec:conclusion}

The canonical problem of time is the block universe in Hamiltonian dress. Read the four-dimensional manifold as the existent, and the constraint correctly reports that nothing evolves; read the evolving three-dimensional world as the existent and its record as the manifold, and the same constraint, deparametrized along the cosmic time that the universe has let us measure, is a true Hamiltonian generating a unitary advance that reproduces the observed flat-$\Lambda$CDM expansion. Lifted beyond the background to the layer's transverse-traceless sector, that same Hamiltonian advances the graviton tower of the closed-$S^3$ layer unitarily---the propagating degree of freedom a fundamental observer reads as the flat-$\Lambda$CDM graviton, with matter the bend of its slicing rather than an ontological content. The equations never changed. What changed was the answer to a single question---what exists---and that answer is not a matter of taste but one the occurrence/existence distinction and the isotropy of the CMB settle together. The problem of time was the canonical symptom of a category error; read without it, the formalism was describing a real cosmic time all along. Read so, the dissolution is one of a family: the same distinction that turns the frozen constraint into a true Hamiltonian is shown, in the companion framework paper's first synthesis, to dissolve a wider set of general relativity's standing problems together---the problem of time among them, and the unitary evolution established here underwriting the resolution of the black-hole information paradox on the realised, globally hyperbolic spacetime~\cite{JanzenCRframework}. That places the canonical face within the framework's one synthesis: the maximally symmetric de~Sitter substrate the framework reads at once as general relativity's solution space~\cite{JanzenRange}, as the discrete $CPT$ and charge structure, and as the gauge sector on its conjugate real form, is read here on its \emph{canonical} constraint---where the deparametrized generator is the layer's true Hamiltonian, and the quantum of action enters the gravitational sector at the seam, scaled by $\Lambda$ alone~\cite{JanzenGeometricCore}.

\begin{thebibliography}{99}

\bibitem{JanzenBoundary} D.~Janzen, \emph{The de Sitter substrate and the Standard Model: a four converging routes to a geometric boundary on what one maximally-symmetric Lorentzian substrate's isometry does and does not force, and the framework and gauge on its two real forms}, companion paper (P13).
\bibitem{JanzenCRframework} D.~Janzen, \emph{Collapsed matter must become a universe: the necessary and sufficient augmentation of general relativity into a description of structural existence, proven for gravitational collapse of any symmetry}, companion paper (P7).

\bibitem{JanzenBHcausality}
D.~Janzen,
\emph{The event horizon is a metric singularity: a missing definition in general relativity}, companion paper (P1).

\bibitem{JanzenAlgebroid}
D.~Janzen,
\emph{General relativity's constraint algebra as the symmetric-space structure of a de Sitter substrate}, companion paper (P12).

\bibitem{JanzenGeometricCore}
D.~Janzen,
\emph{Reached through the imaginary, real by construction: the maximally symmetric substrate as the geometric--ontological core}, companion paper (p0).

\bibitem{JanzenSlicing}
D.~Janzen,
\emph{The de Sitter substrate and its slicing curve: horizons as turning points, and de Sitter and Schwarzschild as two readings of one slicing}, companion paper (P3).

\bibitem{JanzenOperator}
D.~Janzen,
\emph{Covariance of geometries over the de Sitter substrate: the slicing operator, the vacuum kernel, and matter as the bend of the cut}, companion paper (P8).

\bibitem{JanzenRange}
D.~Janzen,
\emph{The range of the de Sitter slicing operator: surjectivity onto the symmetry-reducible sector of general relativity, the Kerr--NUT--(A)dS vacuum kernel, and the wall of inhomogeneity}, companion paper (P9).

\bibitem{JanzenDynamics}
D.~Janzen,
\emph{Why the cut bends: the dynamics of the de Sitter cut, the confined gravitational wave, the generative boundary of the slicing operator, and chirality as the turning of the polarization plane}, companion paper (P11).

\bibitem{JanzenGroupoid}
D.~Janzen,
\emph{The Schwarzschild--de Sitter description groupoid: generators, relations, and Schwarzschild as the asymmetric realisation}, companion paper (P5).

\bibitem{Janzen2025}
D.~Janzen,
\emph{Beyond Space-Time},
(2025).

\bibitem{JanzenWildGoose}
D.~Janzen,
\emph{Albert Einstein's Wild Goose Chase},
cosmiCave.org (2025).

\bibitem{JanzenFortress}
D.~Janzen,
\emph{The Fortress of Misunderstanding Physics},
cosmiCave.org (2025).

\bibitem{ADM1962}
R.~Arnowitt, S.~Deser, and C.~W. Misner,
\emph{The dynamics of general relativity},
in: \emph{Gravitation: An Introduction to Current Research}, ed.~L.~Witten, Wiley, New York (1962) 227--265; arXiv:gr-qc/0405109.

\bibitem{Dirac1964}
P.~A.~M. Dirac,
\emph{Lectures on Quantum Mechanics},
Belfer Graduate School of Science, Yeshiva University, New York (1964).

\bibitem{DeWitt1967}
B.~S. DeWitt,
\emph{Quantum theory of gravity.~I.~The canonical theory},
Phys.~Rev.~\textbf{160} (1967) 1113--1148.

\bibitem{Kuchar1992}
K.~V. Kucha\v{r},
\emph{Time and interpretations of quantum gravity},
in: \emph{Proc.~4th Canadian Conf.~on General Relativity and Relativistic Astrophysics}, eds.~G.~Kunstatter, D.~Vincent, and J.~Williams, World Scientific, Singapore (1992); repr.~Int.~J.~Mod.~Phys.~D~\textbf{20} (2011) 3--86.

\bibitem{Isham1993}
C.~J. Isham,
\emph{Canonical quantum gravity and the problem of time},
in: \emph{Integrable Systems, Quantum Groups, and Quantum Field Theories}, Kluwer, London (1993) 157--287; arXiv:gr-qc/9210011.

\bibitem{BrownKuchar1995}
J.~D. Brown and K.~V. Kucha\v{r},
\emph{Dust as a standard of space and time in canonical quantum gravity},
Phys.~Rev.~D~\textbf{51} (1995) 5600--5629; arXiv:gr-qc/9409001.

\bibitem{PageWootters1983}
D.~N. Page and W.~K. Wootters,
\emph{Evolution without evolution: dynamics described by stationary observables},
Phys.~Rev.~D~\textbf{27} (1983) 2885--2892.

\bibitem{Rovelli2004}
C.~Rovelli,
\emph{Quantum Gravity},
Cambridge University Press, Cambridge (2004).

\bibitem{BarbourEnd}
J.~Barbour,
\emph{The End of Time: The Next Revolution in Physics},
Oxford University Press, Oxford (1999).

\bibitem{Aghanim2020}
Planck Collaboration (N.~Aghanim et~al.),
\emph{Planck 2018 results.~VI.~Cosmological parameters},
Astron.~Astrophys.~\textbf{641} (2020) A6.

\bibitem{JanzenModernParallax}
D.~Janzen,
\emph{The modern parallax: the redshift-isotropy floor, the empirical forcing of cosmic time and uniform expansion, and the measured resolution of the relativity of simultaneity}, companion paper (P4).

\bibitem{JanzenShadowExistence}
D.~Janzen,
\emph{The shadow of existence: scientific theory-choice as an empirically grounded discipline, calibrated on the historical record}, companion paper (P6).

\bibitem{GibbonsHawking1977}
G.~W.~Gibbons and S.~W.~Hawking,
\emph{Cosmological event horizons, thermodynamics, and particle creation},
Phys. Rev. D \textbf{15}, 2738--2751 (1977).

\bibitem{HartleHawking1976}
J.~B.~Hartle and S.~W.~Hawking,
\emph{Path-integral derivation of black-hole radiance},
Phys. Rev. D \textbf{13}, 2188--2203 (1976).

\bibitem{Weyl1910}
H.~Weyl,
\emph{\"Uber gew\"ohnliche Differentialgleichungen mit Singularit\"aten und die zugeh\"origen Entwicklungen willk\"urlicher Funktionen},
Math.~Ann.~\textbf{68} (1910) 220--269.

\bibitem{ReedSimon1975}
M.~Reed and B.~Simon,
\emph{Methods of Modern Mathematical Physics II: Fourier Analysis, Self-Adjointness},
Academic Press, New York (1975).

\bibitem{JanzenCosmogenesis} D.~Janzen, \emph{Cosmogenesis: the Big Bang as a synthesis of Cosmological Relativity's deductively forced consequences, and the primordial light-element abundances it produces}, companion paper (P16).
%  r2376+c54.9: entries below were cited but undefined -- the citations printed as [?]
\bibitem{JanzenCRcosmology} D.~Janzen, \emph{The cosmology of Cosmological Relativity: the expansion history from causal reassignment on a de~Sitter background, the cosmogenesis branch point, and the scalar perturbation sector}, companion paper (P15).
\bibitem{JanzenCircle}
D.~Janzen,
\emph{The Schwarzschild and de Sitter circle: the intrinsic geometry of one homogeneous ring---the event horizon and the curvature singularity at $r=0$ its two poles}, companion paper (P2).
\end{thebibliography}


\clearpage
\section*{Appendix R\quad Computational Receipts}\label{app:receipts}
\addcontentsline{toc}{section}{Appendix R: Computational Receipts}
\small\noindent Each computation marked \rcptmarker\ in the text is verified by a runnable receipt in the \texttt{receipts/} directory that ships with this work; status \textsf{OK} = verified (traced, run, output matches the claim). This appendix is generated from the receipt index, not maintained by hand.\par\medskip
\begingroup\renewcommand{\arraystretch}{1.25}
\begin{description}
\item[\label{rcpt:P10_minisuperspace_friedmann}\texttt{P10\_minisuperspace\_friedmann}]\hfill\textsf{[OK]}\\
\textit{sec:deparam} \ (deparametrized minisuperspace Hphys=(2pi/3)p\_a\textasciicircum{}2/a-(L/8pi)a\textasciicircum{}3 generates Friedmann (adot/a)\textasciicircum{}2=(8pi/3)rho\_d/a\textasciicircum{}3+L/3 via Hamilton's equations). 
\emph{Computes:} adot=4pi p\_a/3a; Hphys autonomous=conserved; eliminate p\_a -> Friedmann; dust a\textasciicircum{}-3; Lambda-sign control
 \emph{Bound:} the true-Hamiltonian structure is non-empty
\ \emph{Run:} \texttt{python3 receipts/P10\_canonical\_time/P10\_minisuperspace\_friedmann.py}
\item[\label{rcpt:P10_graviton_lift}\texttt{P10\_graviton\_lift}]\hfill\textsf{[OK]}\\
\textit{sec:lock} \ (closed-S\textasciicircum{}3 graviton lift: tensor-harmonic action S\_n=1/2 int dT a\textasciicircum{}3[phidot\textasciicircum{}2-(mu\textasciicircum{}2/a\textasciicircum{}2)phi\textasciicircum{}2] Legendre-transforms to Hphys\_n=pi\textasciicircum{}2/2a\textasciicircum{}3+1/2 a mu\textasciicircum{}2 phi\textasciicircum{}2; mode eq phiddot+3(adot/a)phidot+(mu\textasciicircum{}2/a\textasciicircum{}2)phi=0). 
\emph{Computes:} pi=a\textasciicircum{}3 phidot; Legendre H; EL mode eq; Hamilton eqns combine to it; flat-a oscillator control
 \emph{Bound:} unitary cosmic-time evolution of the graviton tower
\ \emph{Run:} \texttt{python3 receipts/P10\_canonical\_time/P10\_graviton\_lift.py}
\item[\label{rcpt:LIFT_adiabatic_correction}\texttt{LIFT\_adiabatic\_correction}]\hfill\textsf{[OK]}\\
\textit{eq:adiabatic-exponent} \ (THE ADIABATIC CORRECTION. Exact Euclidean suppression for mode n is exp(-int omega\_n ds), with omega\_n = mu\_n / a (P10 sec:lock) and a = \textbar{}r\textbar{} along the lift. omega DIVERGES at the branch point (a->0). Does the exponent converge? And is the evolution adiabatic (\textbar{}dw/ds\textbar{}/w\textasciicircum{}2 << 1)?). 
\emph{Computes:} runs clean; registered from its own docstring and a live run
 \emph{Bound:} description is the receipt's own wording, condensed; not independently re-derived here
\ \emph{Run:} \texttt{python3 receipts/P10\_canonical\_time/LIFT\_adiabatic\_correction.py}
\item[\label{rcpt:LIFT_gravitational_action}\texttt{LIFT\_gravitational\_action}]\hfill\textsf{[OK]}\\
\textit{eq:grav-action} \ (THE CALIBRATION. On-shell S = int p\_a da with p\_a = (3/4pi) a adot (from P10's H). Continuing tau = i s : S = -i S\_E with S\_E\textasciicircum{}grav = (3/4pi) int r r'\textasciicircum{}2 ds along the SAME trajectory. Compare with the contour action S\_E\textasciicircum{}c = int [ (1/2) r'\textasciicircum{}2 + f/2 ] ds = 1.489 (alpha=1).). 
\emph{Computes:} runs clean; registered from its own docstring and a live run
 \emph{Bound:} description is the receipt's own wording, condensed; not independently re-derived here
\ \emph{Run:} \texttt{python3 receipts/P10\_canonical\_time/LIFT\_gravitational\_action.py}
\item[\label{rcpt:LIFT_instanton_action}\texttt{LIFT\_instanton\_action}]\hfill\textsf{[OK]}\\
\textit{eq:euclidean-action} \ (FAMILY 10, final item: does the lift have a variational principle of its own? Claim: S\_E = int ds [ (1/2) r'\textasciicircum{}2 + V(r) ], V = f/2, with Euclidean first integral (1/2)r'\textasciicircum{}2 - V = -1/2. Test the closed form r(s) = -(2 M a\textasciicircum{}2)\textasciicircum{}\{1/3\} \textbar{}sin(3s/2a)\textbar{}\textasciicircum{}\{2/3\} against BOTH.). 
\emph{Computes:} runs clean; registered from its own docstring and a live run
 \emph{Bound:} description is the receipt's own wording, condensed; not independently re-derived here
\ \emph{Run:} \texttt{python3 receipts/P10\_canonical\_time/LIFT\_instanton\_action.py}
\item[\label{rcpt:P10_the_straddle_does_not_need_a_floor}\texttt{P10\_the\_straddle\_does\_not\_need\_a\_floor}]\hfill\textsf{[OK]}\\
\textit{sec:lock} \ (**THE STRADDLE DOES NOT NEED A FLOOR, AND THE PARAGRAPH'S TWO STATEMENTS ABOUT Gamma-hat ARE ABOUT DIFFERENT OPERATORS.** The coupled sector promotes the boundary coefficient from the c-number gamma <= 1/4 to an operator Gamma-hat on the tower. **The paragraph asserted it 'bounded below by gamma' where it builds the straddle, and listed 'whether it is bounded below' as OPEN four hundred characters later.** Both are right of different operators: the LEADING-ORDER Gamma-hat = gamma + c sum\_n pi\_n\textasciicircum{}2 is a sum of squares over a positive coefficient, hence spectrum in [gamma, oo); the COMPLETE one carries the cubic and higher self-interactions, which enter at the same inverse-square order and are not sign-definite. ** AND THE REPAIR COSTS THE ARGUMENT NOTHING, WHICH IS THE POINT: the direct-integral decomposition of -d\textasciicircum{}2/dx\textasciicircum{}2 + Gamma-hat/x\textasciicircum{}2 uses only that the spectrum MEETS BOTH SIDES of 3/4 -- limit point at or above, limit circle below -- and a spectral projection onto \{Gamma-hat < 3/4\} exists whether or not the spectrum has a floor.** Exhibited on a two-channel model in which the cubic takes one channel: the union is unbounded below and both spectral subspaces stay non-trivial.). 
\emph{Computes:} the leading-order infimum asserted equal to gamma and the operator asserted monotone in k; the complete model asserted to break the floor and asserted driven down at large momentum; both spectral subspaces asserted non-trivial in both cases
 \emph{Bound:} **bound: the model is a diagonal two-channel stand-in for the singular part of the interacting graviton Hamiltonian, not that Hamiltonian. It establishes what the decomposition NEEDS, not what the spectrum IS -- and the spectrum is exactly the frontier the paragraph names (`L-165`)**
\ \emph{Run:} \texttt{python3 receipts/P10\_canonical\_time/P10\_the\_straddle\_does\_not\_need\_a\_floor.py}
\item[\label{rcpt:P10_the_straddle_is_computed}\texttt{P10\_the\_straddle\_is\_computed}]\hfill\textsf{[ALL PASS]}\\
\textit{sec:lock} \ (THE STRADDLE AS A COMPUTED FACT: spec(Gamma-hat)=[gamma,inf) from the displayed form, and gamma<=1/4<3/4 puts spectrum strictly below the threshold while unboundedness puts it strictly above, so both sides are occupied for every admissible gamma; the 3/4 exhibited by solving the indicial equation rather than quoted). 
\emph{Computes:} r3803
 \emph{Bound:} 59
\ \emph{Run:} \texttt{python3 receipts/P10\_canonical\_time/P10\_the\_straddle\_is\_computed.py}
\item[\label{rcpt:P10_the_adiabatic_residual_at_low_n_is_bounded_by_the_towers_own_floor}\texttt{P10\_the\_adiabatic\_residual\_at\_low\_n\_is\_bounded\_by\_the\_towers\_own\_floor}]\hfill\textsf{[ALL PASS]}\\
\textit{sec:dissolution} \ (**THE ADIABATIC RESIDUAL AT LOW n IS BOUNDED BY THE TOWER'S OWN FLOOR, AND IT REACHES NOTHING OBSERVABLE.** The suppression falls monotonically with harmonic index (2.8e-4 at n=2, 5.9e-6 at n=3, \textasciitilde{}2 orders per step after), so the mixing the adiabatic projection neglects can only carry amplitude into MORE-suppressed modes; and the S\textasciicircum{}3 tensor tower has no mode below n=2, so the least-suppressed mode has no source beneath it. **The bound is therefore one-sided by construction rather than by estimate.** Worst-case transfer of the whole O(eps) fraction out of n=2 leaves that mode MORE suppressed and does not raise the ceiling. Ceiling on transmitted power = unmixed n=2 = **7.9e-8**, and below 1e-3 even on the naive exponent.). 
\emph{Computes:} r4231
 \emph{Bound:} 61
\ \emph{Run:} \texttt{python3 receipts/P10\_canonical\_time/P10\_the\_adiabatic\_residual\_at\_low\_n\_is\_bounded\_by\_the\_towers\_own\_floor.py}
\item[\label{rcpt:P10_gamma_hat_is_bounded_below}\texttt{P10\_gamma\_hat\_is\_bounded\_below}]\hfill\textsf{[OK]}\\
\textit{sec:lock} \ (**THE COMPLETE BOUNDARY OPERATOR IS BOUNDED BELOW, AND THE REASON IS THE DEPARAMETRIZATION ITSELF: THE SCALE FACTOR IS THE SUPERMETRIC'S ONE NEGATIVE DIRECTION.** P10 lists 'the spectrum of Gamma-hat, including whether it is bounded below' among what is open in the interacting tower; that clause is separable and is settled here. **The DeWitt form tr((g pi)\textasciicircum{}2) - lambda (tr g pi)\textasciicircum{}2 has exactly ONE negative eigenvalue on the round metric, realised by the identity direction --- the conformal one, which is the whole indefiniteness of superspace and the reason Wheeler-DeWitt is hyperbolic.** On g-TRACELESS momenta the trace term drops and the form is exactly tr((g pi)\textasciicircum{}2), which for L the Cholesky factor of g equals tr(S\textasciicircum{}2) with S = L\textasciicircum{}T pi L SYMMETRIC --- a sum of squares of real eigenvalues, zero only at pi = 0. **Positive definite for every non-degenerate spatial metric, verified as an IDENTITY on 400 random pairs rather than as a sample.** **So the cubic term's apparent unboundedness is a TRUNCATION ARTEFACT: pi\textasciicircum{}2(1 - lambda phi + ...) is the expansion of pi\textasciicircum{}2/(1 + lambda phi), and the full coefficient is positive wherever the spatial metric is non-degenerate.** The deparametrization that produces a true Hamiltonian is the same move that produces a bounded-below one --- both are the removal of the conformal direction, and P10 argues the first without noticing it had the second.). 
\emph{Computes:} the signature asserted one negative; the Cholesky identity and positivity asserted on 400 valid (g, pi) pairs; min Q asserted positive on three metrics including a 625:1 anisotropy; the truncation asserted to go negative where the full expression does not
 \emph{Bound:} **bound: positivity holds on NON-DEGENERATE spatial metrics. At the degenerate boundary of superspace the inverse metric blows up and the argument says nothing --- which is exactly the a -> 0 boundary, handled instead by the thermal condition. Two different jobs, neither substituting. NOT TOUCHED: the UV definition of the tower sums, and the closed-form nonlinear solution**
\ \emph{Run:} \texttt{python3 receipts/P10\_canonical\_time/P10\_gamma\_hat\_is\_bounded\_below.py}
\item[\label{rcpt:I9_the_adiabatic_parameter_diverges_where_the_action_integral_converges}\texttt{I9\_the\_adiabatic\_parameter\_diverges\_where\_the\_action\_integral\_converges}]\hfill\textsf{[OK]}\\
\textit{sec:lock} \ (I9 --- THE ADIABATIC PARAMETER DIVERGES WHERE THE ACTION INTEGRAL CONVERGES. Integrable-systems field bake, probe I9. On P10's own near-branch-point scaling abs(r) \textasciitilde{} s\textasciicircum{}(2/3), the adiabaticity parameter abs(dw/ds)/w\textasciicircum{}2 goes as s\textasciicircum{}(-1/3) and DIVERGES, so the adiabatic invariant E/omega is not conserved through the branch point; the suppression exponent int omega ds nevertheless converges as S\textasciicircum{}(1/3), because omega \textasciitilde{} s\textasciicircum{}(-2/3) is integrable. Two different powers, two different questions: the finiteness P10 verifies comes from the action integral and not from slow variation, so the correction is semiclassical -- which is what P15 calls the same object. Two adversarial controls, one for each half.). 
\emph{Computes:} runs clean; eight verdicts including controls that fire in opposite directions
 \emph{Bound:} COMPUTES: the scaling abs(r) \textasciitilde{} s\textasciicircum{}(2/3) is P10's own; A and mu carried symbolically, no value pinned
\ \emph{Run:} \texttt{python3 receipts/P10\_canonical\_time/I9\_the\_adiabatic\_parameter\_diverges\_where\_the\_action\_integral\_converges.py}
\item[\label{rcpt:P10_gamma_is_one_quarter_and_is_the_maximum}\texttt{P10\_gamma\_is\_one\_quarter\_and\_is\_the\_maximum}]\hfill\textsf{[OK]}\\
\textit{sec:lock} \ (gamma = 1/4 is the MAXIMUM of the natural ordering family, so (1,1) holds for every ordering). 
\emph{Computes:} reduces T\_s exactly, shows the first-derivative term vanishes, and gets Gamma(s) with its maximum; 11 checks
 \emph{Bound:} NOT-A-PAPER-CLAIM --- verification of a standing claim; node 17 withdrew its own caveat
\ \emph{Run:} \texttt{python3 receipts/P10\_canonical\_time/P10\_gamma\_is\_one\_quarter\_and\_is\_the\_maximum.py}
\item[\label{rcpt:D1_the_boundary_is_per_fibre_and_the_UV_is_over_fibres}\texttt{D1\_the\_boundary\_is\_per\_fibre\_and\_the\_UV\_is\_over\_fibres}]\hfill\textsf{[OK]}\\
\textit{sec:lock} \ (PO-6 re-weighted not closed: the boundary closure is PER-FIBRE and the UV is OVER-FIBRES, and the clause c54.129 answered is the one the argument does not need). 
\emph{Computes:} asserts the direct-integral structure at source, computes the per-fibre deficiency across the threshold, and checks the straddle under an unbounded-below operator (18 checks)
 \emph{Bound:} NOT-A-PAPER-CLAIM --- a bounded re-statement of a PROTECTED row; PO-6 stays open
\ \emph{Run:} \texttt{python3 receipts/L165\_interacting\_tower/D1\_the\_boundary\_is\_per\_fibre\_and\_the\_UV\_is\_over\_fibres.py}
\item[\label{rcpt:D2_the_UV_degree_is_quartic_and_the_IR_is_free}\texttt{D2\_the\_UV\_degree\_is\_quartic\_and\_the\_IR\_is\_free}]\hfill\textsf{[OK]}\\
\textit{sec:lock} \ (PO-6s UV clause measured: the tower sum diverges QUARTICALLY, the ordinary zero-point degree, and compactness buys the IR rather than the UV). 
\emph{Computes:} takes the tower and mode action from P10 at source, derives the linear growth of mu\_n and <pi\_n\textasciicircum{}2>, and tracks the partial sums against N\textasciicircum{}4/4 (13 checks)
 \emph{Bound:} NOT-A-PAPER-CLAIM --- a measurement of one PO-6 clause; PO-6 stays open
\ \emph{Run:} \texttt{python3 receipts/L165\_interacting\_tower/D2\_the\_UV\_degree\_is\_quartic\_and\_the\_IR\_is\_free.py}
\item[\label{rcpt:S1_the_subleading_log_divergence_is_gauss_bonnet_and_the_ledger_survives}\texttt{S1\_the\_subleading\_log\_divergence\_is\_gauss\_bonnet\_and\_the\_ledger\_survives}]\hfill\textsf{[OK]}\\
\textit{---} \ (PO-6 fixed-background half: the SUB-LEADING (log) one-loop divergence on the tower's own de Sitter slicing is the Gauss-Bonnet term (topological in D=4), so it forces NO curvature-squared coupling outside the one-constant ledger --- extending L-809's quartic result to the log order; on de Sitter Weyl\textasciicircum{}2=0, and the on-shell graviton divergence reduces to GB (t Hooft-Veltman/Christensen-Duff)). 
\emph{Computes:} verifies R\textasciicircum{}2/Ric\textasciicircum{}2/Riem\textasciicircum{}2=144/36/24 per alpha\textasciicircum{}4 (S50) and Weyl\textasciicircum{}2=0, that GB=24/alpha\textasciicircum{}4 is the topological Euler density, the on-shell R=4Lambda, and the ledger-survival conclusion (4 checks)
 \emph{Bound:} NOT-A-PAPER-CLAIM --- the fixed-background half of PO-6; extends L-809; does NOT touch the coupled-sector dark half (54's/definitional)
\ \emph{Run:} \texttt{python3 receipts/L816\_po6\_subleading\_survives/S1\_the\_subleading\_log\_divergence\_is\_gauss\_bonnet\_and\_the\_ledger\_survives.py}
\item[\label{rcpt:S1_the_one_constant_ledger_survives_on_the_running_layer_not_only_fixed_de_sitter}\texttt{S1\_the\_one\_constant\_ledger\_survives\_on\_the\_running\_layer\_not\_only\_fixed\_de\_sitter}]\hfill\textsf{[OK]}\\
\textit{sec:lock} \ (PO-6 RUNNING-LAYER half (the residue S50 named: does the one-constant basis survive on a background whose curvature RUNS?): on a(T)=sinh\textasciicircum{}\{2/3\}(3T/2alpha), where R runs (R->12/alpha\textasciicircum{}2 only asymptotically), the pure-graviton log divergence reduces to \{Gauss-Bonnet (topological) + Lambda-renorm + Newton/ell\_P-renorm + field-equation-proportional (field-redef removable)\} with NO irreducible curvature-squared coupling outside \{Lambda,G\} --- for ANY one-loop coefficients. Extends L-816 off fixed de Sitter, via background-independent algebraic identities (not R=const)). 
\emph{Computes:} confirms Weyl\textasciicircum{}2=0 and R runs on the sinh\textasciicircum{}\{2/3\} layer, the topological (53/90)E\_GB split, the Lambda=0 identity (7/20)Ric\textasciicircum{}2+(1/120)R\textasciicircum{}2=E\_mn f\textasciicircum{}mn exactly, the coefficient-independent reduction of a1 Riem\textasciicircum{}2+a2 Ric\textasciicircum{}2+a3 R\textasciicircum{}2 to \{E\_GB,Lam\textasciicircum{}2,LamR,EOM\} with remainder 0, and the vacuum-dS consistency with L-816 (6 checks)
 \emph{Bound:} NOT-A-PAPER-CLAIM --- the running-layer half of PO-6; F5 reserves the verdict; pure graviton (matter as bend), does NOT enter the quantised-background sector or claim one-loop finiteness
\ \emph{Run:} \texttt{python3 receipts/L818\_po6\_running\_layer/S1\_the\_one\_constant\_ledger\_survives\_on\_the\_running\_layer\_not\_only\_fixed\_de\_sitter.py}
\item[\label{rcpt:S1_the_ledgers_one_irreducible_survivor_is_the_conformal_coupling_that_turns_on_at_the_shear}\texttt{S1\_the\_ledgers\_one\_irreducible\_survivor\_is\_the\_conformal\_coupling\_that\_turns\_on\_at\_the\_shear}]\hfill\textsf{[OK]}\\
\textit{sec:lock} \ (PO-6: L-818's boundary LOCATED with a coefficient (answers 56's r2764 routing). The graviton log counterterm = (7/40)Weyl\textasciicircum{}2 + (149/360)GB + (1/8)R\textasciicircum{}2; GB is topological and R\textasciicircum{}2 reduces (L-818), so the ONE irreducible piece is the conformal (7/40)Weyl\textasciicircum{}2 --- it is 0 on the shear-free FRW layer (which is WHY L-818's remainder was exactly 0 and the ledger survived) and turns on at 2nd order in the shear, where Weyl\textasciicircum{}2=4sigma\textasciicircum{}2 (56 S10/r2743). The shear IS the TT-graviton tower = P10's open interacting tower; so L-818's caveat, r2743's shear and P10's open item are ONE boundary, now carrying (7/40)Weyl\textasciicircum{}2). 
\emph{Computes:} decomposes the graviton b\_4 into \{Weyl\textasciicircum{}2,GB,R\textasciicircum{}2\} (Weyl\textasciicircum{}2 coeff 7/40), shows only Weyl\textasciicircum{}2 is neither topological nor ledger-reducible, and computes Weyl\textasciicircum{}2=0 on the isotropic layer and O(shear\textasciicircum{}2) on an axisymmetric anisotropic model (no linear term) (5 checks)
 \emph{Bound:} NOT-A-PAPER-CLAIM --- locates L-818's boundary and names the open coupling; F5 reserves the verdict; does NOT close the interacting tower (it names it); coefficients cited (t Hooft-Veltman)
\ \emph{Run:} \texttt{python3 receipts/L819\_shear\_boundary\_coupling/S1\_the\_ledgers\_one\_irreducible\_survivor\_is\_the\_conformal\_coupling\_that\_turns\_on\_at\_the\_shear.py}
\item[\label{rcpt:S1_the_weyl_magnitude_is_the_graviton_c_anomaly_and_the_parity_odd_partner_is_a_theta_term}\texttt{S1\_the\_weyl\_magnitude\_is\_the\_graviton\_c\_anomaly\_and\_the\_parity\_odd\_partner\_is\_a\_theta\_term}]\hfill\textsf{[OK]}\\
\textit{sec:lock} \ (PO-6 OWED \#472: the VALUE of the one shear counterterm (Weyl\textasciicircum{}2, c54.219) is 7/40 --- the graviton's conformal-anomaly c-coefficient (b\_4 in the \{Weyl\textasciicircum{}2,GB,R\textasciicircum{}2\} basis). And its parity-odd partner, the Pontryagin density R\textasciitilde{}R that a CIRCULAR polarisation reveals (nonzero, handedness-odd; a linear mode gives zero), is a topological theta-term at constant coefficient (Jackiw-Pi C-tensor prop nabla(coeff)=0; no D=4 gravitational anomaly), so it does NOT enter the field equations --- the dynamical count stays ONE). 
\emph{Computes:} computes R\textasciitilde{}R on linear vs circular TT waves by a finite-difference Riemann (0 / +4.5 / -4.5, parity-odd), pins the b\_4 Weyl\textasciicircum{}2 coefficient 7/40, and cites the topological/theta-term status (Jackiw-Pi, Alvarez-Gaume-Witten) (4 checks)
 \emph{Bound:} NOT-A-PAPER-CLAIM --- OWED \#472, one value + one status on the fixed-background shear; the parity-odd loop coefficient is zero for the pure graviton in D=4; interacting tower untouched
\ \emph{Run:} \texttt{python3 receipts/L821\_weyl\_coefficient\_value/S1\_the\_weyl\_magnitude\_is\_the\_graviton\_c\_anomaly\_and\_the\_parity\_odd\_partner\_is\_a\_theta\_term.py}
\item[\label{rcpt:S2_c1_and_c2_are_one_and_the_source_answers_more}\texttt{S2\_c1\_and\_c2\_are\_one\_and\_the\_source\_answers\_more}]\hfill\textsf{[OK]}\\
\textit{sec:euclidean} \ (C1 and C2 are ONE condition joined by so --- the list is six not seven --- and the same P10 passage states that the free TT sector is a harmonic oscillator whose Hamiltonian is bounded below, which PO-6 asks for). 
\emph{Computes:} asserts the joined sentence verbatim, confirms PO-6s row asks whether it is bounded below and does not carry the answer, and quotes the confinement clause (11 checks)
 \emph{Bound:} NOT-A-PAPER-CLAIM --- corrects a count and surfaces a partial answer; joint satisfiability untested
\ \emph{Run:} \texttt{python3 receipts/L165\_defining\_the\_sum/S2\_c1\_and\_c2\_are\_one\_and\_the\_source\_answers\_more.py}
\item[\label{rcpt:S3_c6_is_a_theorem_not_a_condition}\texttt{S3\_c6\_is\_a\_theorem\_not\_a\_condition}]\hfill\textsf{[OK]}\\
\textit{sec:lock} \ (the C6/C7 tension cannot arise: the per-fibre structure is DERIVED from a commutation relation, not a requirement on a measure --- so the conditions list is five plus one theorem). 
\emph{Computes:} asserts D1s commutation, direct-integral and fibre-by-fibre statements verbatim, and S1s misclassification of C6 as a condition (8 checks)
 \emph{Bound:} NOT-A-PAPER-CLAIM --- reclassifies a list entry; joint satisfiability untested
\ \emph{Run:} \texttt{python3 receipts/L165\_defining\_the\_sum/S3\_c6\_is\_a\_theorem\_not\_a\_condition.py}
\item[\label{rcpt:S4_the_open_half_is_the_floor}\texttt{S4\_the\_open\_half\_is\_the\_floor}]\hfill\textsf{[OK]}\\
\textit{sec:lock} \ (PO-6s open half is the FLOOR: the direct-integral decomposition needs only that both sides of the threshold are occupied, and P10 leaves whether the complete Gamma-hat is bounded below open in its own voice). 
\emph{Computes:} asserts P10s uses-only clause, its leaves-open clause, the sum-of-squares form and the structural commutation verbatim, and the PO-6 rows own wording (7 checks)
 \emph{Bound:} NOT-A-PAPER-CLAIM --- narrows a register row and corrects r2611s caveat
\ \emph{Run:} \texttt{python3 receipts/L165\_defining\_the\_sum/S4\_the\_open\_half\_is\_the\_floor.py}
\item[\label{rcpt:D3_the_floor_does_not_survive_the_cubic}\texttt{D3\_the\_floor\_does\_not\_survive\_the\_cubic}]\hfill\textsf{[OK]}\\
\textit{sec:lock} \ (PO-6s floor question answered NO at the order the paper names: the cubic enters as pi\textasciicircum{}2 phi, a square times a SIGNED field, so Gamma-hat is unbounded below wherever phi < -c/g3). 
\emph{Computes:} asserts P10s leaves-open clause, its leading-order form and its naming of the cubic verbatim, and computes the sign structure and divergence (8 checks)
 \emph{Bound:} NOT-A-PAPER-CLAIM --- answers a narrowed dark half; the state question is untouched
\ \emph{Run:} \texttt{python3 receipts/L165\_interacting\_tower/D3\_the\_floor\_does\_not\_survive\_the\_cubic.py}
\item[\label{rcpt:D4_the_state_fixer_has_its_own_threshold}\texttt{D4\_the\_state\_fixer\_has\_its\_own\_threshold}]\hfill\textsf{[OK]}\\
\textit{sec:lock} \ (the state-fixing mechanism has its own threshold: nu = sqrt(Gamma + 1/4) is real only for Gamma >= -1/4, below which both exponents go complex and no regular branch exists --- and r2651s operator reaches -47.8). 
\emph{Computes:} asserts P10s thermal-regularity sentence and its common-to-every-fibre clause verbatim, computes nu across Gammas range, and recomputes r2651s minimum (11 checks)
 \emph{Bound:} NOT-A-PAPER-CLAIM --- sharpens PO-6s successor question; whether the region is realised is untouched
\ \emph{Run:} \texttt{python3 receipts/L165\_interacting\_tower/D4\_the\_state\_fixer\_has\_its\_own\_threshold.py}
\item[\label{rcpt:D50_the_floor_does_survive_and_the_paper_said_so_232_revisions_earlier}\texttt{D50\_the\_floor\_does\_survive\_and\_the\_paper\_said\_so\_232\_revisions\_earlier}]\hfill\textsf{[OK]}\\
\textit{sec:deparam} \ ( **`D3` (r2651) HAS THE FLOOR BACKWARDS --- P10 CALLS THE CUBIC'S UNBOUNDEDNESS AN "ARTEFACT OF TRUNCATION" IN THE SAME PARAGRAPH, AND HAS SINCE r2419.** The cubic is the first-order term of \$\textbackslash\{\}pi\textasciicircum{}2/(1+\textbackslash\{\}lambda\textbackslash\{\}phi)\$, whose full coefficient is positive on non-degenerate metrics --- so the floor survives and r2652's sub-\$(-1/4)\$ question has no object. Lead `L-542`; VEIN `L-165`). 
\emph{Computes:} verifies termwise to eighth order that the cubic IS the expansion's first-order term; samples both forms on 200,000 identical points (truncation \$-892.6\$, resummed never below \$\textbackslash\{\}gamma\$); six source checks in P10 including the affirmative verdict and its interior-only SCOPE; dates the clause by reading r2419's own copy of the file; and re-runs `P10\_gamma\_hat\_is\_bounded\_below`, which passes
 \emph{Bound:} LATENT --- no new physics; the answer was in the paper. NOT claimed: that `PO-6` closes (the ultraviolet definition of the tower sums is open in the same sentence), or that the floor is proved for the full interacting theory.  `D3`/`D4` and the `PO-6` row are NOT touched --- 56's band and a protected row; routed
\ \emph{Run:} \texttt{python3 receipts/L165\_interacting\_tower/D50\_the\_floor\_does\_survive\_and\_the\_paper\_said\_so\_232\_revisions\_earlier.py}
\item[\label{rcpt:S50_the_counterterm_basis_is_one_dimensional_because_the_background_family_is}\texttt{S50\_the\_counterterm\_basis\_is\_one\_dimensional\_because\_the\_background\_family\_is}]\hfill\textsf{[OK]}\\
\textit{sec:deparam} \ (**`PO-6`'s UV HALF: THE COUNTERTERM BASIS IS ONE-DIMENSIONAL BECAUSE THE ADMITTED BACKGROUND FAMILY IS.** Every curvature invariant on a maximally symmetric background is a pure power of \$1/\textbackslash\{\}alpha\textasciicircum{}2\$, and the substrates form a one-parameter family, so a divergence of any degree needs exactly one counterterm --- which the framework carries.  Scoped: the tower lives on the layer, whose \$R\$ runs. Lead `L-543`; VEIN `L-165`). 
\emph{Computes:}  **AMENDED c54.211 (`L-544`): c54.210's SCOPE IS WITHDRAWN.** It put P15's observable rate where P10's slicing belongs; P10's own \$a(T)=\textbackslash\{\}alpha\textbackslash\{\}cosh(T/\textbackslash\{\}alpha)\$ has \$R=12/\textbackslash\{\}alpha\textasciicircum{}2\$ CONSTANT, so the degeneracy holds on the tower's own background and the result is STRONGER than claimed. The real limit is the coupled sector, where the scale factor is quantized and there is no fixed background. Both geometries computed side by side so the error stays checkable. --- computes the quadratic invariants at \$D=4,5,6\$ and pins their \$\textbackslash\{\}alpha\$-exponent at exactly 4; five standing-sentence checks (p0's invariant fact, p0's foliation, `L-533`'s scale-only descent); computes the \$\textbackslash\{\}sinh\textasciicircum{}\{2/3\}\$ layer's Ricci scalar and asserts it NON-constant with a \$12H\textasciicircum{}2\$ late limit; and counts `counterterm` in P10 from git HEAD (0) against now (7)
 \emph{Bound:} LATENT, and an ASSEMBLY not a discovery. NOT claimed: that `PO-6` closes, that the sub-leading divergences vanish, that this is renormalisability, or that the layer breaks it --- **that is the question**
\ \emph{Run:} \texttt{python3 receipts/L165\_defining\_the\_sum/S50\_the\_counterterm\_basis\_is\_one\_dimensional\_because\_the\_background\_family\_is.py}
\item[\label{rcpt:S7_the_successor_was_misaimed}\texttt{S7\_the\_successor\_was\_misaimed}]\hfill\textsf{[OK]}\\
\textit{sec:tower} \ (L-543 INHERITS r2677s WITHDRAWN OBJECT: it asks about a RUNNING background and the free tower sits on P10s own slicing at R = 12/alpha\textasciicircum{}2, CONSTANT --- and the layers late limit IS that value, so it asymptotes to it). 
\emph{Computes:} computes both Ricci scalars, shows the layers late limit equals P10s constant under H -> 1/alpha, and pins the successors wording and the back-reaction clause from the row (6 checks)
 \emph{Bound:} NOT-A-PAPER-CLAIM --- withdraws a successor this line registered; the back-reaction question stands
\ \emph{Run:} \texttt{python3 receipts/L165\_interacting\_tower/S7\_the\_successor\_was\_misaimed.py}
\item[\label{rcpt:S8_the_floor_follows}\texttt{S8\_the\_floor\_follows}]\hfill\textsf{[OK]}\\
\textit{sec:lock} \ (PO-6 HAS BEEN CARRYING THE WRONG QUESTION: the complete Gamma-hat inherits its floor from the same non-degeneracy r2671 established, so boundedness is answered YES --- what remains is WHERE the spectrum sits relative to 3/4). 
\emph{Computes:} asserts P10s promotion, inverse-square and Weyl-alternative clauses verbatim, quotes the straddle receipts does-not-need-a-lower-bound line, and shows the truncated coefficient goes negative where the full one stays positive (7 checks)
 \emph{Bound:} NOT-A-PAPER-CLAIM --- the spectrum is untouched; P10 explicitly declines the floor and the decomposition survives without it
\ \emph{Run:} \texttt{python3 receipts/L165\_interacting\_tower/S8\_the\_floor\_follows.py}
\item[\label{rcpt:S9_the_ordering_decides_it}\texttt{S9\_the\_ordering\_decides\_it}]\hfill\textsf{[OK]}\\
\textit{sec:lock} \ (PO-6s SPECTRAL QUESTION IS DECIDED BY OPERATOR ORDERING, which P10 never names: normal-ordered puts the vacuum at 1/4 (BELOW 3/4, freedom survives), symmetric puts N=1 exactly AT 3/4 --- and 3/4 = 1/4 + 1/2, the free coefficient plus one zero-point quantum). 
\emph{Computes:} computes both ordering branches against the threshold, verifies the 1/4+1/2 identity, counts normal-order and symmetric-order at zero in the paper, and pins the both-sides-occupied clause (9 checks)
 \emph{Bound:} NOT-A-PAPER-CLAIM --- neither ordering is asserted correct; P10s decomposition survives either
\ \emph{Run:} \texttt{python3 receipts/L165\_interacting\_tower/S9\_the\_ordering\_decides\_it.py}
\item[\label{rcpt:Q1_the_degeneracy_is_conformal_flatness_not_maximal_symmetry_so_no_scale_factor_can_break_it}\texttt{Q1\_the\_degeneracy\_is\_conformal\_flatness\_not\_maximal\_symmetry\_so\_no\_scale\_factor\_can\_break\_it}]\hfill\textsf{[OK]}\\
\textit{sec:tower} \ (PO-6's DARK HALF HAS AN OBJECT: r2677 stated TWO degeneracies as one. The quadratic trio collapses because the deficit is exactly Weyl\textasciicircum{}2 and EVERY FRW is conformally flat for EVERY a(T) -- not because of maximal symmetry -- so no scale factor can break it and back-reaction cannot; the different-dimension half does fail. The real limit is the SHEAR). 
\emph{Computes:} R, Ric\textasciicircum{}2 and Riem\textasciicircum{}2 computed from the FRW metric with a(T) a free function at k = +1, 0, -1 (Weyl\textasciicircum{}2 identically zero); the Gauss-Bonnet density shown an exact total derivative, sqrt(g)(R\textasciicircum{}2-4Ric\textasciicircum{}2+Riem\textasciicircum{}2) = d/dT[24(a'\textasciicircum{}3/3+a')]; the two relations solved to give int Ric\textasciicircum{}2 = A/3 - chi/2 and int Riem\textasciicircum{}2 = A/3 - chi; S50's D=4 row reproduced from the metric (144/36/24); P15's running layer recomputed from its own metric with the deficit zero by a second route in u = exp(3Ht/2); and Bianchi I giving Weyl\textasciicircum{}2 = 4 sigma\textasciicircum{}2 + O(sigma\textasciicircum{}4). CONTROLS: SdS gives Weyl\textasciicircum{}2 = 48M\textasciicircum{}2/r\textasciicircum{}6 so the identity is not vacuous; the layer's R\textasciicircum{}2 is t-dependent so the different-dimension degeneracy really does fail; and the numeric residual is a PRECISION-SCALING test (3.06e-37 at 40 digits, 2.24e-77 at 80) rather than a fitted threshold
 \emph{Bound:} not that the coupled sector is renormalizable -- this says which functionals a divergence can NEED, not that they are absorbable; no heat-kernel coefficient computed; Bianchi I is a HOMOGENEOUS shear, fixing the ORDER at which conformal flatness fails and not the mode-by-mode statement on the tower; and not a closure -- PO-6 stays open
\ \emph{Run:} \texttt{python3 receipts/L549\_coupled\_counterterms/Q1\_the\_degeneracy\_is\_conformal\_flatness\_not\_maximal\_symmetry\_so\_no\_scale\_factor\_can\_break\_it.py}
\item[\label{rcpt:S10_the_halves_meet_at_the_shear}\texttt{S10\_the\_halves\_meet\_at\_the\_shear}]\hfill\textsf{[OK]}\\
\textit{sec:lock} \ (PO-6s TWO DECLARED HALVES MEET AT THE SHEAR: the counterterm half ends where conformal flatness ends (C\textasciicircum{}2 = 4 sigma\textasciicircum{}2 + O(sigma\textasciicircum{}4)) and P10 names the tower as the TRANSVERSE-TRACELESS graviton modes --- which ARE the shear). 
\emph{Computes:} checks C\textasciicircum{}2 vanishes at sigma=0, is quadratic-leading and non-zero away from it, confirms the source is trace-free, and quotes P10s tower identification twice (6 checks)
 \emph{Bound:} NOT-A-PAPER-CLAIM --- the second-order basis is not computed; cc54 derived the coefficient and its structure is checked here
\ \emph{Run:} \texttt{python3 receipts/L165\_interacting\_tower/S10\_the\_halves\_meet\_at\_the\_shear.py}
\item[\label{rcpt:S11_the_ordering_question_dissolves}\texttt{S11\_the\_ordering\_question\_dissolves}]\hfill\textsf{[OK]}\\
\textit{sec:lock} \ (PO-6s ORDERING QUESTION DISSOLVES: P10 gives deficiency indices (1,1) INDEPENDENTLY of ordering with the coefficient <= 1/4 across the natural family, closes the residual via the horizons Hartle-Hawking state, and answers the coupled straddle fibre by fibre). 
\emph{Computes:} pins all six of P10s own statements --- the independence, the family bound, the non-uniqueness, the Friedrichs selection, the fibre-by-fibre condition, and what remains open (6 checks)
 \emph{Bound:} NOT-A-PAPER-CLAIM --- the interacting tower is still open; the correction is to WHICH problem the row owes
\ \emph{Run:} \texttt{python3 receipts/L165\_interacting\_tower/S11\_the\_ordering\_question\_dissolves.py}
\item[\label{rcpt:P10_ordering_selection_is_external}\texttt{P10\_ordering\_selection\_is\_external}]\hfill\textsf{[OK]}\\
\textit{sec:lock / P7 frontier:quantum} \ (**PO-15 THE ORDERING EXHAUSTION: nothing internal to the construction selects the operator ordering (min Gamma-hat = 1/4 normal vs 3/4 symmetric), and the reason is that selecting it would be solving the cosmological-constant problem.** Completes S11 (boundary CLOSURE is ordering-independent) and the r3014 thermal-state elimination (the thermal state acts downstream). ** THE PIVOT: ** the ordering is BULK-INERT (normal vs symmetric differ by the c-number zero-point 1/2 hbar omega\_n, a global phase under the deparametrized evolution) and BOUNDARY-PHYSICAL (the difference surfaces only in the horizon coefficient, 1/4 -> 3/4 = 1/4 + 1/2, exactly one zero-point quantum), so "which ordering" IS "does the tower's zero-point energy gravitate at the horizon". ** THE EXHAUSTION: ** (1) the thermal state and (3) the seam's characteristic structure act DOWNSTREAM (they close the extension GIVEN Gamma-hat; a regular branch x\textasciicircum{}(1/2+nu) exists for BOTH orderings); (2) the substrate isometry, (5) positivity, (6) the single-scale ledger and (7) covariance under configuration redefinition each respect BOTH orderings; (4) the deparametrization makes normal ordering AVAILABLE (a preferred vacuum) but not mandatory, and by SOLVING the constraint removes the anomaly-freedom lever a Wheeler-DeWitt quantization would have used. The two external naturalness heuristics DISAGREE (continuity-with-free -> 1/4, predictivity -> 3/4). ** RESULT: ** the choice is EXTERNAL because it IS the cc problem reached from inside -- an epistemic gap (one datum: does vacuum energy gravitate), NOT an ontological family, so admissible under the epistemic reading.). 
\emph{Computes:} the ordering difference asserted = the zero-point c-number 1/2 hbar omega; the c-number-is-a-global-phase fact asserted numerically (identical amplitude moduli to 1e-12, ratio = e\textasciicircum{}-ict); the boundary shift 1/4 -> 3/4 asserted to be exactly one zero-point quantum with 3/4 the threshold; the regular branch nu=sqrt(Gamma+1/4) asserted real for BOTH orderings (thermal closure ordering-independent); both floors asserted finite; the seven-candidate enumeration stated with a structural disqualification each and no survivor
 \emph{Bound:} PO-15 answered as the row asked -- the choice is genuinely external, the enumeration is the evidence, the externality is the cc problem localized. NO ordering is picked and none is claimed derivable; the interacting-tower definition (PO-6) stays separately open; builds on S11/r3014, not a re-derivation of the Friedrichs selection
\ \emph{Run:} \texttt{python3 receipts/P10\_canonical\_time/P10\_ordering\_selection\_is\_external.py}
\item[\label{rcpt:S12_the_boundaries_coincide}\texttt{S12\_the\_boundaries\_coincide}]\hfill\textsf{[OK]}\\
\textit{sec:lock} \ (cc54s L-818 BOUNDARY, r2743s SHEAR FINDING AND P10s OPEN ITEM ARE ONE BOUNDARY: an FRW layer is shear-free by construction, the degeneracy ends at the shear, and P10 names the transverse-traceless tower as what remains open). 
\emph{Computes:} verifies cc54s Step 1 (Weyl\textasciicircum{}2 = 0 with R running, all three k) and pins P10s two namings of the tower (5 checks)
 \emph{Bound:} NOT-A-PAPER-CLAIM --- cc54s five sympy steps are not rerun; the coincidence is by identification, each link a corpus statement
\ \emph{Run:} \texttt{python3 receipts/L165\_interacting\_tower/S12\_the\_boundaries\_coincide.py}
\item[\label{rcpt:S13_the_ledger_does_not_reach_second_order}\texttt{S13\_the\_ledger\_does\_not\_reach\_second\_order}]\hfill\textsf{[OK]}\\
\textit{sec:lock} \ (AT SECOND ORDER IN THE SHEAR THE QUADRATIC BASIS IS GENUINELY THREE-DIMENSIONAL: GB and C\textasciicircum{}2 span rank 2 of 3, C\textasciicircum{}2 vanishes identically shear-free so L-818 never had to route it, and it has NONE of L-818s three routes). 
\emph{Computes:} computes the rank and nullspace of \{GB, C\textasciicircum{}2\}, confirms C\textasciicircum{}2 = 0 at all three k on the shear-free layer, and pins p0s one-constant claim as a statement about the FACES (6 checks)
 \emph{Bound:} BOUNDED NEGATIVE --- one order of one expansion; not a verdict on the construction, and the coefficient is not computed
\ \emph{Run:} \texttt{python3 receipts/L165\_interacting\_tower/S13\_the\_ledger\_does\_not\_reach\_second\_order.py}
\item[\label{rcpt:S1_the_shear_needs_exactly_one_new_counterterm_and_my_own_count_was_one_too_many}\texttt{S1\_the\_shear\_needs\_exactly\_one\_new\_counterterm\_and\_my\_own\_count\_was\_one\_too\_many}]\hfill\textsf{[OK]}\\
\textit{sec:tower} \ (PO-6's OWED SHEAR CALCULATION RUN: on a propagating TT mode C\textasciicircum{}2 is O(h\textasciicircum{}2) and carries FOUR derivatives, but delta\textasciicircum{}(1)R = 0 makes int sqrt(g) R\textasciicircum{}2 a shift of Einstein-Hilbert -- so the shear costs exactly ONE new counterterm, Weyl-squared, and c54.215's working count of two was one too many). 
\emph{Computes:} R, Ric\textasciicircum{}2, Riem\textasciicircum{}2 and C\textasciicircum{}2 built from the corpus's own L801/N1 TT-wave-on-de-Sitter metric and expanded to O(eps\textasciicircum{}2); the derivative content fixed by freezing the oscillatory factors and taking the total degree in (H,k,omega), giving 4 against sigma\textasciicircum{}2's 2; delta\textasciicircum{}(1)R shown to vanish and det g shown eps-independent, hence R\textasciicircum{}2 at second order equals 2 Rbar times R at second order, pointwise; the span solved from Gauss-Bonnet plus the Weyl definition, giving B and C both in span of A and C\textasciicircum{}2; and the Pontryagin density computed for linear, circular and opposite-handed polarisation by a finite-difference pipeline cross-checked against a symbolic evaluation to six digits. CONTROLS: the linear mode returns zero for the parity-odd invariant, which is what catches the incomplete basis; the handedness flips its sign; and at eps = 0 everything reduces to c54.215's one dimension
 \emph{Bound:} no heat-kernel COEFFICIENT is computed -- this says which functionals a divergence can need, not with what coefficient; the mode is a plane-wave TT on flat-sliced de Sitter and not P10's S\textasciicircum{}3 harmonics, so the COUNT is mode-independent while the COEFFICIENT is not; the Pontryagin term's total-derivative status is cited, not computed; its consequences for the corpus's chirality result are not raised; and not a closure -- PO-6 stays open
\ \emph{Run:} \texttt{python3 receipts/L553\_the\_shear\_counterterm/S1\_the\_shear\_needs\_exactly\_one\_new\_counterterm\_and\_my\_own\_count\_was\_one\_too\_many.py}
\item[\label{rcpt:D1_the_degeneracy_carrying_the_quartic_was_never_derived_and_its_constant_is_the_component_count}\texttt{D1\_the\_degeneracy\_carrying\_the\_quartic\_was\_never\_derived\_and\_its\_constant\_is\_the\_component\_count}]\hfill\textsf{[OK]}\\
\textit{sec:tower} \ (THE DEGENERACY CARRYING PO-6's QUARTIC WAS NEVER DERIVED: asserted in P10 and asserted-unchecked in D2. Derived from Peter-Weyl on the parallelizable S\textasciicircum{}3 it is g(n) = 2(n-1)(n+3), n >= 2, so the shell contribution is 2n\textasciicircum{}3 -- and the constant is the propagating-component count, which makes D2's "any tensor rank" false). 
\emph{Computes:} the level-j totals of frame-valued fields on S\textasciicircum{}3 = SU(2) for frame-spin 0, 1, 2, coming out as 1, 3, 5 times (2j+1)\textasciicircum{}2 -- the component counts, which is the check the decomposition had to pass; the TT part as the two extreme summands, 2(2j+1)\textasciicircum{}2, verified to be exactly 2/5 of the symmetric-tracefree total; the closed form by eigenvalue with its floor g(2) = 10 matching P10's independently-asserted n >= 2; the cumulative sum and its leading coefficient 2/3; and Weyl's law calibrated on the scalar tower, whose degeneracy follows exactly from Peter-Weyl, with the ratio running 1.01505 to 1.00038. CONTROLS: the scalar tower's leading coefficient 1/3 against the tensor tower's 2/3, which is what refutes "any tensor rank"; and a REPORTED NEGATIVE -- matching Weyl to subleading order fails to discriminate the candidate closed forms, because on a closed curved manifold that term carries the curvature through the heat kernel
 \emph{Bound:} not that the eigenvalue assignment is re-derived, mu\_n\textasciicircum{}2 = n(n+2)-2 being P10's and taken as given; no heat-kernel coefficient, still none computed anywhere; this is the FREE spectrum's mode count and not the interacting tower; and not a closure -- PO-6 stays open, this replaces an unchecked input with a derived one
\ \emph{Run:} \texttt{python3 receipts/L554\_tower\_degeneracy/D1\_the\_degeneracy\_carrying\_the\_quartic\_was\_never\_derived\_and\_its\_constant\_is\_the\_component\_count.py}
\end{description}\endgroup


\end{document}
